\documentclass[11pt]{article}
\usepackage[T1]{fontenc}
\usepackage{lmodern}
\usepackage{amsmath,amssymb,amsthm,mathtools}
\usepackage{mathrsfs}
\usepackage[margin=1in]{geometry}
\usepackage[colorlinks=true,linkcolor=blue,citecolor=blue,urlcolor=blue]{hyperref}
\usepackage{enumitem}
\allowdisplaybreaks
\setlength{\emergencystretch}{2em}
\numberwithin{equation}{section}
\newtheorem{theorem}{Theorem}[section]
\newtheorem{proposition}[theorem]{Proposition}
\newtheorem{corollary}[theorem]{Corollary}
\newtheorem{lemma}[theorem]{Lemma}
\theoremstyle{definition}
\newtheorem{definition}[theorem]{Definition}
\theoremstyle{remark}
\newtheorem{remark}[theorem]{Remark}
\DeclareMathOperator{\Rea}{Re}
\DeclareMathOperator{\Ima}{Im}
\DeclareMathOperator{\Log}{Log}
\newcommand{\RH}{\mathrm{RH}}
\newcommand{\dd}{\,d}
\title{Collision Exclusion and Collective Attraction\protect\\
in the Riemann Heat Flow:\protect\\
Reductions, Quantitative Tools, and Endpoint Obstructions}
\author{Drafted for Edward Baker}
\date{9 October 2026}
\begin{document}
\maketitle
\begin{abstract}
We assemble analytic reductions and proved comparison tools for a proposed
descent of the de Bruijn--Newman bound. A positive threshold forces a finite
real multiple zero, while positive-kernel and theta-cutoff examples delimit
generic exclusion mechanisms. A holomorphic normalized approximation gives
explicit disk, reflection, cutoff, and derivative errors. In the shrinking
sector $\kappa=t\log(x/(4\pi))\in[1,2]$, the complete analytic error decays
exponentially. Genuine phase cancellation removes growing terminal blocks,
with a larger removal for the derivative coordinate. Weighted reflected
frequency estimates give positive mean square for the complete derivative
on intervals of length $Dt^2e^{2\mathfrak a/t}$, where
$\mathfrak a=\kappa(4-\kappa)/16$. This interval remains too long for
the available local curvature argument. Independent-phase and complete
multiplicative-twist models admit joint zeros; their scope is distinguished
from the genuine height orbit. At an all-real multiple zero, deflation and
the forced mirror root give a necessary fourth-derivative inequality.
Its full-sum and value-core approximation errors vanish, supplying an
additional conditional signed exclusion test. A polynomial heat flow
saturates that inequality at an arbitrary positive threshold. Separately,
a sharp exact attraction-field comparison and uniform all-zero counts
give global field and landing tools, with worked effective and compact
examples. The opposite signed inequality, higher-multiplicity control,
and coverage outside the shrinking sector remain open. No proof of RH,
descent to zero, new numerical record, or literature priority is claimed.
\end{abstract}
\tableofcontents

\section{Setup and scope of the reductions}
\label{sec:setup}
Use the standard normalization
\begin{equation}\label{eq:heat}
 H_t(z)=\int_0^\infty e^{tu^2}\Phi(u)\cos(zu)\dd u,
 \qquad H_0(z)=\frac18\xi\left(\frac12+\frac{iz}{2}\right),
 \qquad \partial_tH_t=-H_t'',
\end{equation}
where
\begin{equation}\label{eq:theta}
 \Phi(u)=\sum_{n\ge1}\phi_n(u),\qquad
 \phi_n(u)=(2\pi^2n^4e^{9u}-3\pi n^2e^{5u})
                    e^{-\pi n^2e^{4u}}.
\end{equation}
The full kernel extends to an even function, by the theta functional
equation, and is positive on the real axis. The representation, threshold
theorem, and quantitative approximation used below are imported from the
classical theory and \cite{polymath,newmanwu}. The integral and all its
finite parameter derivatives converge locally uniformly: the decay in
\eqref{eq:theta} dominates every fixed Gaussian and exponential weight.
Thus the flow is jointly entire in the complex variables $t,z$, and is real
and even in $z$ at real times. In particular,
\begin{equation}\label{eq:imaginarypositive}
 H_t(iy)=\int_0^\infty e^{tu^2}\Phi(u)\cosh(yu)\dd u>0
 \quad(t,y\in\mathbb R).
\end{equation}

There is a finite constant $\Lambda$ such that all zeros of $H_t$ are real
exactly for $t\ge\Lambda$. RH is equivalent to $\Lambda\le0$, and
Rodgers--Tao prove $\Lambda\ge0$ \cite{rodgerstao}. De Bruijn's strip
contraction, in this normalization, says that a strip of half-width $a$
at time $t_0$ becomes a strip of half-width
\begin{equation}\label{eq:strip}
 \sqrt{\max\{a^2-2h,0\}}\qquad\text{at time }t_0+h,
 \quad h\ge0.
\end{equation}
The original critical strip gives $a=1$ at time zero and hence
$\Lambda\le1/2$. A zero-free half-plane $\Rea s>\theta$, together with
the functional equation, gives $a=2\theta-1$ when $1/2\le\theta\le1$.
Repeated use of \eqref{eq:strip} alone preserves the absolute endpoint
$t_0+a^2/2$.

The results proved here are deductions from these imported theorems.
Their novelty in the literature has not been established. The
projected-field framework of \cite{planat} is credited in
Section~\ref{sec:field}; its numerical certificates and reported endpoints
are not inputs. Our central distinction is between a strictly improved
positive-time landing estimate and a mechanism excluding every positive
collision. The latter would be an RH-strength theorem.

\section{Finite collisions and the endpoint obstruction}
\label{sec:collisions}
The uniform localization input needed here is
\cite[Theorem~1.5(i)]{polymath}: for an absolute $C$, a zero
$H_t(x+iy)=0$ with $0<t\le1/2$ and $x\ge e^{C/t}$ has $y=0$.
Evenness handles negative $x$. Uniformity for $t\ge\varepsilon>0$ is
essential; mere finiteness of the exceptional set separately at each
time would not give the following compactness argument.

\begin{proposition}[A positive threshold entails a finite collision]
\label{prop:finitecollision}
If $\Lambda>0$, there is $x_*\in\mathbb R\setminus\{0\}$ such that
\begin{equation}\label{eq:collision}
 H_\Lambda(x_*)=H_\Lambda'(x_*)=0.
\end{equation}
\end{proposition}
\begin{proof}
Choose $\Lambda/2\le t_j<\Lambda$ increasing to $\Lambda$, and a
nonreal zero $z_j$ at each time. Strip contraction confines their imaginary
parts to $[-1,1]$, while the imported localization confines their real
parts to $[-e^{2C/\Lambda},e^{2C/\Lambda}]$. A subsequence converges
to $z_*$, and continuity gives $H_\Lambda(z_*)=0$. The threshold property
makes $z_*=x_*$ real. If the root were simple, the implicit-function
theorem would give a unique nearby zero for real times near $\Lambda$;
conjugation would force it to be real. This contradicts the nearby
nonreal $z_j$. Finally \eqref{eq:imaginarypositive} excludes $x_*=0$.
\end{proof}

For later use, the imported local splitting theorem
\cite[Proposition~3.1]{polymath} states that a real zero of multiplicity
$m\ge2$ at $(T,x_*)$ has nearby roots
\begin{equation}\label{eq:hermite}
 x_*+\sqrt{2(t-T)}\,\lambda_j+O(|t-T|),
\end{equation}
where the $\lambda_j$ are the distinct real roots of the probabilists'
Hermite polynomial $\operatorname{He}_m$. For $t<T$, at least one of
these roots is nonreal. For $t>T$, all are real and simple sufficiently
close to $T$.

\begin{corollary}[Exact endpoint target]\label{cor:rhcollision}
RH is equivalent to the absence of a common real zero of $H_t,H_t'$ for
every $t>0$. If the time-zero zeros lie in $|\Ima z|\le a$, it suffices
to exclude such a zero for $0<t\le a^2/2$.
\end{corollary}
\begin{proof}
Under RH, all zeros are real for $t\ge0$. A multiple root at positive
time would contradict this through \eqref{eq:hermite} at a slightly
earlier nonnegative time. If RH fails, the threshold equivalence gives
$\Lambda>0$ and Proposition~\ref{prop:finitecollision} applies.
After $a^2/2$ all roots are real by \eqref{eq:strip}; a later multiple
root is excluded by the same local splitting argument.
\end{proof}

\begin{proposition}[Ordinary collision geometry]\label{prop:double}
At an exact double real zero of a real analytic heat flow at $(T,x_*)$,
the nearby conjugate roots for $t<T$ have the form
\begin{equation}\label{eq:double}
 a(t)\pm iy(t),\qquad a(t)=x_*+O(T-t),\qquad
 y(t)^2=2(T-t)+O((T-t)^2).
\end{equation}
The collision is already transverse as an intersection of the two real
equations $H=H'=0$:
\begin{equation}\label{eq:jacobian}
 \det\frac{\partial(H,H')}{\partial(x,t)}(T,x_*)=(H''(T,x_*))^2>0.
\end{equation}
\end{proposition}
\begin{proof}
Translate $x_*$ to zero and apply real Weierstrass preparation:
$H_t(z)=A(t,z)(z^2+b(t)z+c(t))$, with $A(T,0)\ne0$ and
$b(T)=c(T)=0$. The heat equation at this point gives $c'(T)=-2$.
Completing the square gives
$(z+b(t)/2)^2=2(t-T)+O((t-T)^2)$, proving \eqref{eq:double}.
For \eqref{eq:jacobian}, the two rows of the Jacobian, in the stated
column order $(x,t)$, are $(H',-H'')$ and $(H'',-H''')$.
\end{proof}

Consequently $t+y(t)^2/2=T+O((T-t)^2)$. Finite additional attraction
can improve the quadratic term while remaining compatible with $T>0$.
Transversality, or a theorem that every collision is double, would not
exclude this event. A bound on $\Lambda$ also cannot be propagated
backward into a new strip at an earlier time: \eqref{eq:strip} evolves
forward. These facts are the endpoint obstruction to interpreting a
sequence of strictly improved positive bounds as descent to zero.

\section{A positive-kernel counterexample to a general mechanism}
\label{sec:counterexample}
This section concerns a different entire function, not zeta. It excludes
an argument based only on positive even kernels, rapid decay, a narrow
zero strip, and finitely many exceptional zeros.

\begin{proposition}\label{prop:kernel}
Put $\varphi(u)=e^{-\cosh u}$ and
$q(u)=\varphi^{(4)}(u)+324\varphi(u)$. Then $q$ is even, strictly
positive, strictly decreasing for $u>0$, and double-exponentially
decaying. The flow
\begin{equation}\label{eq:qflow}
 J_t(z)=\int_{\mathbb R}e^{tu^2}q(u)e^{izu}\dd u
\end{equation}
has a finite threshold $0<\Lambda_q\le9/2$. At time zero exactly four
zeros are nonreal, namely $\pm3\pm3i$, and all other zeros are real.
\end{proposition}
\begin{proof}
Write $c=\cosh u\ge1$. Direct differentiation gives
\begin{align*}
 q(u)&=(c^4-6c^3+5c^2+5c+321)e^{-c},\\
 c^4-6c^3+5c^2+5c+321
 &=\frac{(2c-9)^2(4c^2+12c+27)}{16}
   +5(c^2-1)+5(c-1)+\frac{3109}{16}.
\end{align*}
Thus positivity, evenness, and decay hold. Moreover
\[
 q'(u)=-\sinh(u)e^{-c}r(c),\qquad
 r(c)=c^2(c-5)^2-2c^2-5c+316.
\]
For $1\le c\le8$, $r(c)\ge316-128-40=148$; for $c\ge8$,
$r(c)\ge7c^2-5c+316>0$. This proves monotonicity. The defining
integral and all its derivatives converge for every fixed Gaussian
weight. In particular the general de Bruijn kernel assumptions, as
recorded in \cite[Theorem~7]{newmanwu}, are satisfied.

Let $B(z)=\int_{\mathbb R}e^{-\cosh u}e^{izu}\dd u=2K_{iz}(1)$.
The all-real-zero property of this Bessel function is classical; see
\cite{gasper}. Here is a spectral proof of the required special case.
If $K_{iz}(1)=0$, the nonzero decaying function $v(x)=K_{iz}(e^x)$,
$x\ge0$, satisfies
\[
 -v''+e^{2x}v=z^2v,\qquad v(0)=0.
\]
Integration against $\overline v$, with the vanishing boundary term,
gives
\[
 z^2\int_0^\infty|v|^2\dd x
 =\int_0^\infty(|v'|^2+e^{2x}|v|^2)\dd x>0.
\]
The standard decaying large-argument asymptotic of $K_\nu$ shows that
$v$ is not identically zero and justifies both integrals. Thus $z^2$
is positive real and $z$ is real. Four integrations by parts give
\[
 J_0(z)=(z^4+324)B(z).
\]
The four polynomial roots $\pm3\pm3i$ are simple and $B$ has no
nonreal root, giving the asserted zero set and a strip of half-width
three.

De Bruijn's theorem makes the zeros all real at $t\ge9/2$ and preserves
this property upward in time. The all-real property is closed under
locally uniform convergence, by Hurwitz's theorem applied to disks
off the real axis; the limit cannot be identically zero because
$J_t(0)>0$. Its set of times is therefore a closed upper ray.
The four simple nonreal roots persist for small positive $t$, so the
lower endpoint of that ray is strictly positive.
\end{proof}

\begin{corollary}\label{cor:rescale}
For $a>0$, the kernel $q_a(u)=a q(au)$ has
\[
 J_t^{[a]}(z)=J_{t/a^2}(z/a),\qquad
 \Lambda_{q_a}=a^2\Lambda_q>0.
\]
Its time-zero zeros lie in $|\Ima z|\le3a$, and only four are nonreal.
\end{corollary}
\begin{proof}
Substitute $v=au$ in the integral; zeros and the threshold scale as
displayed.
\end{proof}
The scale changes the function. It does not produce successively sharper
strips for one fixed zeta function, and it does not preserve zeta's prime
coefficients, Euler product, explicit formula, or zero density.
Vanishing of all odd endpoint jets also holds for this smooth even
counterexample; that property alone cannot be a zeta-specific sign
mechanism.

\section{The exact theta interface and a cutoff obstruction}
\label{sec:theta}
Define the collision vector and Wronskian on the real axis by
\[
 C_t(x)=H_t(x)^2+H_t'(x)^2,\qquad
 W_t(x)=H_t'(x)^2-H_t(x)H_t''(x).
\]
$C_t(x)>0$ is equivalent to absence of a collision at that point.
$W_t(x)>0$ is sufficient, whereas $W_t(x)\ge0$ allows a common zero.
With $f_t(u)=e^{tu^2}\Phi(u)$, symmetrization and the product-to-sum
identities give the exact formula
\begin{equation}\label{eq:wronskianintegral}
 W_t(x)=\frac14\int_0^\infty\!\int_0^\infty f_t(u)f_t(v)
 \big((u+v)^2\cos(x(u-v))+(u-v)^2\cos(x(u+v))\big)\dd u\dd v.
\end{equation}
The actual squared-integer coefficients from \eqref{eq:theta} are
positive, but the cosines need not be. Kernel positivity therefore does
not give the desired sign term by term.

Another exact interface, obtained by $v=\pi n^2e^{4u}$ in each summand,
is
\begin{equation}\label{eq:exactarithmetic}
 H_t(x)=\frac{\pi^{-1/4}}4\int_\pi^\infty (2v-3)v^{1/4}e^{-v}
 \sum_{n\le\sqrt{v/\pi}}n^{-1/2}
 e^{t\ell_n(v)^2/16}\cos(x\ell_n(v)/4)\dd v,
 \quad \ell_n(v)=\log\frac{v}{\pi n^2}.
\end{equation}
For the $j$th spatial derivative, multiply the summand by
$(\ell_n(v)/4)^j$ and replace its cosine by
$\cos(x\ell_n(v)/4+j\pi/2)$. Absolute convergence justifies these
operations. The finite sum depends on $v$ and the outer integral is
infinite. Formula~\eqref{eq:exactarithmetic} is not an Euler product for
$H_t$, and a family estimate for these weights would be an additional
theorem.

\begin{proposition}[Every fixed theta cutoff has a negative tail]
\label{prop:thetacutoff}
For $N\ge1$, set $\Phi_N=\sum_{n\le N}\phi_n$ and
$H_{t,N}(x)=\int_0^\infty e^{tu^2}\Phi_N(u)\cos(xu)\dd u$.
Then, uniformly for $t$ in any compact real interval,
\begin{align}\label{eq:thetaasymp}
 H_{t,N}(x)&=-A_Nx^{-2}+O(x^{-4}),&
 H_{t,N}'(x)&=2A_Nx^{-3}+O(x^{-5}),\\
 H_{t,N}''(x)&=-6A_Nx^{-4}+O(x^{-6}),&
 W_{t,N}(x)&=-2A_N^2x^{-6}+O(x^{-8}),\nonumber
\end{align}
where, with $a_n=\pi n^2$,
\begin{equation}\label{eq:endpointjet}
 A_N=\Phi_N'(0)
 =\sum_{n>N}a_n(8a_n^2-30a_n+15)e^{-a_n}>0.
\end{equation}
In particular $W_{t,N}(x)<0$ for all sufficiently large $|x|$.
\end{proposition}
\begin{proof}
Direct differentiation gives
$\phi_n'(0)=-a_n(8a_n^2-30a_n+15)e^{-a_n}$.
The full even kernel has $\Phi'(0)=0$, so Gaussian convergence of the
derivative sum gives \eqref{eq:endpointjet}. Its tail terms are positive
because $n>N\ge1$ implies $a_n\ge4\pi$, beyond the larger root
$(15+\sqrt{105})/8$ of the quadratic. For
$f(u)=e^{tu^2}\Phi_N(u)$, all derivatives are integrable uniformly
on compact time intervals and $f'(0)=A_N$. Repeated integration by
parts gives
\[
 \int_0^\infty f(u)\cos(xu)\dd u
 =-\frac{f'(0)}{x^2}+\frac{f'''(0)}{x^4}+O(x^{-6}).
\]
Apply integration by parts separately to $-u f(u)\sin(xu)$ and
$-u^2 f(u)\cos(xu)$. Their first nonzero coefficients are
$2f'(0)$ and $-6f'(0)$, respectively. This avoids differentiating
an uncontrolled asymptotic remainder. Substitution gives the
Wronskian assertion.
\end{proof}

All odd jets of the full kernel, including after multiplication by
$e^{tu^2}$, cancel exactly. The first omitted tail jet is $-A_N$,
so \eqref{eq:thetaasymp} describes a cutoff defect rather than a sign
change of the full zeta Wronskian. Indeed $H_{t,N}$ cannot have all
zeros real at any finite time. Its integral gives growth
$\log\max_{|z|\le r}|H_{t,N}(z)|=O(r\log(r+2))$, hence order at
most one. A real entire function of order less than two with all zeros
real satisfies the Laguerre inequality $F'^2-FF''\ge0$, by its
Hadamard product. This contradicts \eqref{eq:thetaasymp}.

Simple analytic symmetrization is not a remedy at positive time:
$\Phi_N(-u)\sim-3\pi(\sum_{n\le N}n^2)e^{-5u}$ as $u\to\infty$,
so $(\Phi_N(u)+\Phi_N(-u))/2$ is not absolutely heat-integrable
after multiplication by $e^{tu^2}$ for $t>0$. A suitable replacement
must preserve both the cancellations and heat integrability.

\section{A normalized holomorphic collision criterion}
\label{sec:normalization}
We now use a different finite approximation, the effective
Riemann--Siegel representation from \cite[Theorem~1.3]{polymath}.
Its scalar theorem applies to $0<t\le1/2$, $x\ge200$,
$0\le y\le1$. We keep this domain, rewrite the reflected sum
holomorphically, freeze its integer cutoff, and derive the spatial
derivative error by Cauchy's formula.

\subsection{The fixed-cutoff analytic representation}
Put $s=(1-iz)/2$ and use principal logarithms in
\begin{align}\label{eq:m0}
 m_0(s)={}&\Log s+\Log(s-1)-\frac s2\log\pi
 +\log\frac{\sqrt{2\pi}}{16}
 +\left(\frac s2-\frac12\right)\Log\frac s2-\frac s2,\\
 \alpha(s)={}&\frac1{2s}+\frac1{s-1}+\frac12\Log\frac{s}{2\pi},
 \qquad m_t(s)=m_0(s)+\frac t4\alpha(s)^2,
 \qquad M_t(s)=e^{m_t(s)}.\nonumber
\end{align}
For positive real part of $z$, both $s$ and $1-s$ avoid the cut
$(-\infty,1]$. For a fixed $N\ge1$, define
\begin{align}\label{eq:p}
 p_{t,n}(s)&=\exp\left(\frac t4\log^2n
                  -(s+\tfrac t2\alpha(s))\log n\right),&
 P_{t,N}(s)&=\sum_{n\le N}p_{t,n}(s),\\
 G_{t,N}(z)&=M_t(s)P_{t,N}(s)+M_t(1-s)P_{t,N}(1-s).\nonumber
\end{align}
The source theorem uses the natural cutoff
\begin{equation}\label{eq:naturalcutoff}
 n(t,x)=\left\lfloor\sqrt{\frac{x}{4\pi}+\frac t{16}}\right\rfloor.
\end{equation}
Its seemingly conjugated reflected sum equals the second term of
\eqref{eq:p}: $\overline s-y=1-s$, and its correction
$\kappa=\tfrac t2(\alpha(1-s)-\alpha(\overline s))$ converts the
exponent $-\overline{s+t\alpha(s)/2}-\kappa+y$ to
$-(1-s)-t\alpha(1-s)/2$. Thus no conjugate variable remains.
$G_{t,N}$ is holomorphic when $N$ is fixed and has real symmetry.

Define the nonvanishing analytic normalizer and its phase by
\begin{align}\label{eq:symmetricnormalization}
 A_t(z)&=\exp\frac{m_t(s)+m_t(1-s)}2,&
 \theta_t(z)&=\frac{m_t(s)-m_t(1-s)}{2i},\\
 Q_t(z)&=\frac{H_t(z)}{A_t(z)},&
 F_{t,N}(z)&=\frac{G_{t,N}(z)}{A_t(z)}
 =e^{i\theta_t}P_{t,N}(s)+e^{-i\theta_t}P_{t,N}(1-s).\nonumber
\end{align}
On the real axis $A_t(x)>0$ and $\theta_t,Q_t,F_{t,N}$ are real.
The remainder has real symmetry throughout its neighborhood. Since
$A_t$ has no zeros,
\begin{equation}\label{eq:collisionnormalized}
 H_t(x)=H_t'(x)=0\quad\Longleftrightarrow\quad Q_t(x)=Q_t'(x)=0.
\end{equation}
Normalization removes the exponentially small common archimedean factor;
it does not provide a lower bound for the collision vector.

For exact derivative evaluation put
$a_{t,n}(s)=M_t(s)p_{t,n}(s)$ and
$\lambda_{t,n}(s)=(\alpha(s)-\log n)(1+t\alpha'(s)/2)$. Then
\begin{align}\label{eq:normalizedderivative}
 G_{t,N}'&=\frac i2\sum_{n\le N}
  \big(\lambda_{t,n}(1-s)a_{t,n}(1-s)
       -\lambda_{t,n}(s)a_{t,n}(s)\big),\\
 b_t:=\frac{A_t'}{A_t}&=-\frac i4\big(m_t'(s)-m_t'(1-s)\big),
 \qquad F_{t,N}'=G_{t,N}'/A_t-b_tF_{t,N},\nonumber\\
 m_t'(s)&=\alpha(s)(1+t\alpha'(s)/2).\nonumber
\end{align}
The fixed $N$ in these formulas is essential: differentiating a formula
with the jumping cutoff \eqref{eq:naturalcutoff} is invalid at a jump.

\subsection{Explicit disk, reflection, and cutoff errors}
Fix $0<\varepsilon\le\tau\le1/2$, $0<R\le1$, and $X-R\ge200$.
Set
\[
 x_-=X-R,\quad x_+=X+R,\quad q_\pm=\frac{x_\pm}{4\pi},
 \quad V=\frac{x_-}{2},\quad S=\frac12\sqrt{x_+^2+(1+R)^2},
\]
and define
\begin{align}\label{eq:diskconstants}
 \alpha_-&=\frac12\log q_- -V^{-2}>0,&
 A_*&=\frac{3}{2V}+\frac12
           \sqrt{\log^2(S/(2\pi))+\pi^2/4},\\
 D_*&=\frac{3}{2V^2}+\frac1{2V},&
 L_*&=A_*(1+\tau D_*/2),\qquad K_*=e^{RL_*/2}.\nonumber
\end{align}
Throughout the upper enclosing rectangle
$x_-\le\Rea z\le x_+$, $0\le\Ima z\le R$, both $s,1-s$ have
real parts in $[0,1]$, imaginary modulus at least $V$, and modulus
at most $S$. Consequently
\begin{equation}\label{eq:alphabounds}
 \Rea\alpha\ge\alpha_-,\quad |\alpha|\le A_*,\quad
 |\alpha'|\le D_*,\quad |m_t'|\le L_*.
\end{equation}
For the first inequality use
$\Rea(1/(2s))\ge0$, $\Rea(1/(s-1))\ge-V^{-2}$, and $|s|\ge V$;
the others follow directly from \eqref{eq:m0}, including
$|\arg s|\le\pi/2$. Integrating vertically from the real axis gives
\begin{equation}\label{eq:normalizercost}
 |M_t(s)/A_t|,\ |M_t(1-s)/A_t|\le K_*.
\end{equation}

The permitted natural cutoffs are among
\[
 m_-:=\left\lfloor\sqrt{q_-+\varepsilon/16}\right\rfloor,
 \qquad m_+:=\left\lfloor\sqrt{q_++\tau/16}\right\rfloor.
\]
Here $m_-\ge3$ and $m_+-m_-\le1$: the range under the square
root has length at most $1/(2\pi)+1/32<1$, and its square-root
range also has length less than one. Define positive majorants
\begin{align}\label{eq:termmajorants}
 d_n&=\tfrac14\log^2n-\tfrac12\alpha_-\log n,&
 V_n&=n^{-1/2}e^{\max\{\varepsilon d_n,\tau d_n\}},&
 P_n&=n^{R/2}V_n,\\
 \ell_n&=\max\{|\log(q_-/n^2)|,|\log(q_+/n^2)|\},&
 k_*&=\frac{\tau R}{2(x_--6)},&
 U_n&=\exp\frac{\tau^2\ell_n^2/16+0.626}{x_--6.66}-1.\nonumber
\end{align}
The endpoint maximum in $V_n$ preserves the sign of the heat exponent;
its logarithm is affine in $t$. For each permitted $m$, put
\begin{align}\label{eq:eta}
 E_m^{AB}&=\sum_{n\le m}(1+m^{k_*}n^R)V_nU_n,\\
 E_m^C&=q_-^{-1/4}\exp\left(-\frac\varepsilon{16}\log^2q_-
 +\frac{1.24(3^R+3^{-R})}{m-0.125}
 +\frac{3\sqrt{\log^2q_++\pi^2/4}+10.44}{x_--12}\right),\nonumber\\
 \eta_N&=K_*\max_{m_-\le m\le m_+}
 \left(E_m^{AB}+E_m^C
 +2\sum_{n=\min(m,N)+1}^{\max(m,N)}P_n\right).\nonumber
\end{align}
An empty sum is zero. The decimal constants retain their exact meanings
in the imported inequalities of \cite{polymath}; these are not empirical
fit constants.

\begin{proposition}[Holomorphic remainder and derivative control]
\label{prop:diskerror}
For every $t\in[\varepsilon,\tau]$ and $|z-X|\le R$,
\begin{equation}\label{eq:diskerror}
 |Q_t(z)-F_{t,N}(z)|\le\eta_N.
\end{equation}
For real $|x-X|\le a<R$, this also gives
\begin{equation}\label{eq:derivativeerror}
 |Q_t'(x)-F_{t,N}'(x)|\le\frac{\eta_N}{R-a}.
\end{equation}
\end{proposition}
\begin{proof}
On the upper rectangle, apply \cite[Theorem~1.3, equations~(23)--(24)]{polymath}
with its actual cutoff $m=n(t,x)$. The source error is bounded, in
its $M_t(s)$ normalization, by $E_m^{AB}+E_m^C$:
$|\gamma|\le e^{0.02y}(x/(4\pi))^{-y/2}\le1$,
$m^{|\kappa|}\le m^{k_*}$, $n^y\le n^R$, and
$|p_{t,n}(s)|\le V_n$ follow from the source inequalities and
\eqref{eq:alphabounds}. The remaining terms are majorized
monotonically as in \eqref{eq:eta}. Changing to the symmetric
normalizer costs $K_*$ by \eqref{eq:normalizercost}. Changing $m$ to
the fixed $N$ adds or removes exactly the finite terms shown in
\eqref{eq:eta}; each reflected normalized term is at most $K_*P_n$,
since both real parts are at least $(1-R)/2$. Real symmetry of
$Q_t-F_{t,N}$ reflects the upper-half estimate to the lower half-disk.
This proves \eqref{eq:diskerror} on the full disk. Cauchy's estimate
on the radius-$R-a$ circle about an inner real point proves
\eqref{eq:derivativeerror}; the remainder is holomorphic and the
normalizer has no zeros in the neighborhood.
\end{proof}

\begin{corollary}[Cell criterion]\label{cor:cell}
On $t\in[\varepsilon,\tau]$, $|x-X|\le a<R$, either uniform inequality
\[
 |F_{t,N}(x)|>\eta_N\qquad\text{or}\qquad
 |F_{t,N}'(x)|>\eta_N/(R-a)
\]
excludes collisions. A cover may use a different successful alternative
in each cell.
\end{corollary}
\begin{proof}
At a common zero of $Q_t,Q_t'$, both opposite weak inequalities
would follow from Proposition~\ref{prop:diskerror}. Then use
\eqref{eq:collisionnormalized}.
\end{proof}
The uniform cell bounds require interval enclosures or analytic
variation bounds, not point samples. Reflection after an asymmetric
normalization would not preserve the same error; the factor $K_*$
is part of the proof.

\subsection{A leading-pair exclusion and its remaining gap}
For $N=1$, the entire finite tail is still paid in \eqref{eq:eta}, while
\begin{equation}\label{eq:leadingpair}
 F_{t,1}(x)=2\cos\theta_t(x),\qquad
 F_{t,1}'(x)=-2\theta_t'(x)\sin\theta_t(x),\qquad
 \theta_t'(x)=-\tfrac12\Rea m_t'((1-ix)/2).
\end{equation}
Put $\omega_*=(\alpha_- -\tau A_*D_*/2)/2$. If $\omega_*>0$,
then $|\theta_t'|\ge\omega_*$ by \eqref{eq:alphabounds}.
A candidate collision therefore contradicts $\cos^2\theta+\sin^2\theta=1$
whenever
\begin{equation}\label{eq:paircriterion}
 \left(\frac{\eta_1}{2}\right)^2+
 \left(\frac{\eta_1}{2(R-a)\omega_*}\right)^2<1.
\end{equation}

\begin{corollary}[Explicit eventual exclusion]\label{cor:eventual}
For $0<\varepsilon\le1/2$,
\begin{equation}\label{eq:expheight}
 \varepsilon\le t\le\tfrac12,\qquad |x|\ge e^{64/\varepsilon}
 \quad\Longrightarrow\quad(H_t(x),H_t'(x))\ne(0,0).
\end{equation}
\end{corollary}
\begin{proof}
By evenness take $X=x>0$, $\tau=1/2$,
$L=\log(X/(4\pi))$, $R=1/L$, and $a=0$. From $X\ge e^{128}$,
$\log(4\pi)<2.6$, and \eqref{eq:diskconstants}--\eqref{eq:termmajorants},
the following relaxed elementary bounds hold:
\begin{gather*}
 L>125,\quad\log q_->125,\quad\varepsilon\log q_->62,
 \quad R<1/125,\\
 \alpha_-\ge\tfrac12\log q_- -0.001,
 \quad\log m_+\le\tfrac12\log q_-+0.01,\\
 A_*\le0.51L,\quad D_*<0.001,\quad K_*<1.3,
 \quad R\omega_*>0.24.
\end{gather*}
For example $|\log q_\pm-L|\le2R/X$, and
$\log m_+\le\tfrac12\log(q_++1/32)$ give the middle estimates.
Also $L_*/L<0.52$ gives $K_*<e^{0.26}<1.3$;
$\alpha_-/L>0.499$ and $\tau A_*D_*/(2L)<0.000128$ give the
phase estimate. In particular $d_n\le0$ for $n\le m_+$. With
\[
 p=\frac{1-R}{2}+\frac\varepsilon2\alpha_-
                       -\frac\varepsilon4\log m_+>8,
\]
we have $P_n\le n^{-8}$ and hence
\[
 \sum_{n=2}^{m_+}P_n\le\frac1{256}+\frac1{896}.
\]
The bound $\ell_n\le\log X$ and monotonic decrease of
$(\log X)^2/X$ for $X\ge e^{128}$ give $U_n<10^{-40}$.
Since $m_+^{k_*}<2$ and $p-R>7$, the two power sums in
$E_m^{AB}$ give $E_m^{AB}<4\cdot10^{-40}$.
In $E_m^C$ the negative heat exponent is below $-484$ and the
positive terms together are below three, so $E_m^C<e^{-481}<10^{-100}$.
Uniformly over the permitted cutoffs,
\[
 \eta_1<1.3\left(4\cdot10^{-40}+10^{-100}
               +2\left(\frac1{256}+\frac1{896}\right)\right)<0.014.
\]
The left side of \eqref{eq:paircriterion} is at most
$(0.014/2)^2+(0.014/0.48)^2<0.001$. This excludes the collision
at the arbitrary point $X$.
\end{proof}

This corollary rederives known positive-time eventual simplicity with
a deliberately large explicit height; it is not a new bound near zero
or a record. The remaining finite height grows without bound as
$\varepsilon\downarrow0$. No fixed finite certificate removes that
quantifier gap. The effective finite approximation here differs from
the theta cutoff of Section~\ref{sec:theta}: every omitted finite and
analytic term is paid as an error.

\section{The exact attraction field and a sharp probe comparison}
\label{sec:field}
Let $z=x+iy$, $y>0$, be a simple nonreal zero whose imaginary height
is globally maximal, and put $w=y^2$. Zeros are counted with multiplicity.
Canonical-product sums below are grouped under the real and even
symmetries; this is the convention in the imported zero-motion theorem
\cite[Proposition~3.1]{polymath}. Define
\begin{equation}\label{eq:fields}
 E_t(z)=\frac1y\sum_{\rho\ne z,\overline z}
 \frac{y-\Ima\rho}{(x-\Rea\rho)^2+(y-\Ima\rho)^2},
 \qquad
 G_t(z)=\sum_{\rho\ne z,\overline z}
 \frac1{(x-\Rea\rho)^2+4w}.
\end{equation}
$E$ is the exact external field; $G$ is its projected counterpart.
The latter is the field used in the collective-contraction framework
of \cite[Section~2]{planat}. The present comparison below concerns $E$.

\begin{lemma}[Zero motion and projected floor]\label{lem:motion}
At such a zero,
\begin{equation}\label{eq:motion}
 w'=-2-4wE_t(z),\qquad E_t(z)\ge G_t(z)\ge0.
\end{equation}
\end{lemma}
\begin{proof}
Implicit differentiation of $H_t(z(t))=0$ and the canonical product
give $z'=H_t''(z)/H_t'(z)=2\sum_{\rho\ne z}(z-\rho)^{-1}$ in
the grouped convention. The own conjugate contributes $-2$ to $w'$.
The remaining imaginary parts give the definition of $E$.
A real external root at horizontal distance $d$ contributes
$1/(d^2+w)$ to $E$, which exceeds its contribution to $G$.
For an external conjugate pair at heights $\pm v$, $0<v\le y$,
put $s=d^2/y^2$ and $u=1-v^2/y^2\in[0,1]$. Its exact contribution
is $2(s+u)/(wD_E)$, where
$D_E=s^2+2(2-u)s+u^2$, while its projected contribution is
$2/(w(s+4))$. The difference has nonnegative numerator
$(s+u)(s+4)-D_E=3us+u(4-u)$. Summation proves the assertion.
The excluded configuration $s=u=0$ would be the target pair itself.
\end{proof}

For $\eta>y$ the point $x+i\eta$ lies above every zero. Define
\begin{equation}\label{eq:probe}
 L_\eta(t,x)=-\Ima\frac{H_t'}{H_t}(x+i\eta),\qquad
 Q_\eta=\frac{L_\eta}{\eta}-\frac2{\eta^2-w}\ge0,
 \qquad K(c)=\min\{1,(c^2-1)/4\}.
\end{equation}
The subtraction removes exactly the own conjugate pair in the
canonical-product logarithmic derivative.

\begin{theorem}[Sharp exact-field comparison]\label{thm:probe}
Every $\eta>y$ satisfies
\begin{equation}\label{eq:sharpbridge}
 E_t(z)\ge K(\eta/y)
 \left(\frac{L_\eta(t,x)}{\eta}-\frac2{\eta^2-w}\right).
\end{equation}
The factor is sharp for the allowed pairwise geometry. In particular,
if $y<\eta\le\sqrt5\,y$, then
\begin{equation}\label{eq:nearprobe}
 w'\le-(\eta^2-w)\frac{L_\eta(t,x)}{\eta}.
\end{equation}
\end{theorem}
\begin{proof}
For an external pair put $s=d^2/y^2$, $u=1-v^2/y^2$, and
$q=\eta^2/y^2-1>0$. Apart from the common factor $2/y^2$,
the exact and probe contributions are respectively
\[
 e=\frac{s+u}{D_E},\qquad p=\frac{s+q+u}{D_Q},\quad
 D_E=s^2+2(2-u)s+u^2,\quad
 D_Q=s^2+2(q+2-u)s+(q+u)^2.
\]
For $0<q\le4$, the cleared numerator of $4e-qp$ is
\begin{align*}
 &4(s+u)D_Q-q(s+q+u)D_E\\
 &=(4-q)s^3+[12+4(1-u)+q(4-q)+uq]s^2\\
 &\quad+u[12+4(1-u)+12q+2q^2+qu]s
       +u(u+q)[(4-q)u+4q]\ge0.
\end{align*}
For $q\ge4$, the numerator of $e-p$ is
\[
 (s+u)D_Q-(s+q+u)D_E
 =q[s^2+(q-4+6u)s+u(q+u)]\ge0.
\]
A real root has exact-to-probe ratio $(s+q+1)/(s+1)\ge1$.
Summing these termwise comparisons proves \eqref{eq:sharpbridge}.
At $u=0$, the pair ratio tends to $q/4$ as $s\downarrow0$ and
to one as $s\to\infty$, establishing pairwise sharpness. This
does not assert that these extremal configurations occur for zeta.
For $q\le4$, insert the comparison into \eqref{eq:motion}; the
own-pair terms cancel exactly and give \eqref{eq:nearprobe}.
\end{proof}

At $\eta\ge\sqrt5\,y$ the comparison is a consequence of the
existing projected-field bridge \cite[Proposition~2.3]{planat}.
The nearer-probe assertion concerns $E$ and has the explicit smaller
factor; we do not assert it for $G$. An adaptive probe $\eta(w)$ is
evaluated instantaneously and is not differentiated in the zero-motion
formula. For a fixed $\eta>\sqrt{w_0}$ and a uniform bound
$L_\eta\ge L_0$ throughout a region containing every maximizer,
\begin{equation}\label{eq:fixedprobefloor}
 E_t(z)\ge K(\eta/\sqrt{w_0})
 \max\left\{0,\frac{L_0}{\eta}-\frac2{\eta^2-w_0}\right\}
 \quad(0<w\le w_0).
\end{equation}
Both factors on the right of \eqref{eq:sharpbridge} can only improve
as $w$ decreases under these bounds. Spatial and temporal coverage
is part of the hypothesis; a high-height sector alone cannot imply
a global landing endpoint.

\section{Uniform counts, global coverage, and landing}
\label{sec:landing}
Let $N_t(R)$ count all zeros with $0<\Rea\rho\le R$, including
nonreal zeros and multiplicity. The imported
\cite[Theorem~1.5(iv)]{polymath} supplies an absolute $A\ge0$ such that
\begin{equation}\label{eq:count}
 |N_t(R)-p_t(R)|\le A\log(2+R),\qquad 0<t\le\tfrac12,\quad R\ge4\pi,
\end{equation}
where
\[
 p_t(R)=\frac R{4\pi}\log\frac R{4\pi}-\frac R{4\pi}
          +\frac{11}{8}+\frac t{16}\log\frac R{4\pi}.
\]
The error constant is independent of $t$. Endpoint conventions can
be absorbed into $A$ using the same theorem's local count bound.
No numerical value of this constant is supplied here.

\begin{proposition}[A global field from all-zero counts]\label{prop:countfloor}
Set
\begin{equation}\label{eq:countparameters}
 D=8\pi(4A+1),\quad X_* =\max\{(4\pi)^2,2D,3\},
 \quad n=\log X_*,\quad B=X_*+2D.
\end{equation}
For every simple globally maximal nonreal zero at $0<t\le1/2$,
\begin{equation}\label{eq:globalfloor}
 E_t(z)\ge G_t(z)\ge\frac{n}{B^2+4w}.
\end{equation}
\end{proposition}
\begin{proof}
Evenness permits $x\ge0$. If $x\ge X_*$, use only the zeros with
real parts in $(x+D,x+2D]$. They exclude both own-pair roots. On
this interval,
\[
 p_t'(R)=\frac1{4\pi}\log\frac R{4\pi}+\frac{t}{16R}
 \ge\frac{\log x}{8\pi},
\]
since $x\ge(4\pi)^2$. Also $x\ge2D,3$ gives
$2+x+2D\le2+2x\le x^2$, so each endpoint error logarithm is at
most $2\log x$. The count in the block is at least
\[
 \left(\frac{D}{8\pi}-4A\right)\log x=\log x.
\]
Every horizontal distance is at most $2D$, giving
$G_t\ge\log x/(4D^2+4w)$. If $0\le x\le X_*$, use instead
the fixed block $(X_*+D,X_*+2D]$, which has at least $n$ zeros
by the same argument and distances at most $B$. For large $x$,
use $\log x\ge n$ and $B\ge2D$ to obtain the same bound.
Finally apply Lemma~\ref{lem:motion}.
\end{proof}

This proposition uses all-zero counts; replacing them by a count
of real zeros would lose its uniformity near time zero. No horizontal
confinement assumption is needed. Another available floor comes
from the reflected external pair $-x\pm iy$: by
\eqref{eq:imaginarypositive}, $x\ne0$, and
\begin{equation}\label{eq:mirrorfloor}
 E_t(z)\ge G_t(z)\ge\frac1{2(x^2+w)}.
\end{equation}
The density, reflected-pair, and probe fields ordinarily overlap in
their zero contributions. Their floors may be combined by a maximum.
Adding them requires disjoint blocks or a separately proved joint
comparison.

\begin{theorem}[State-dependent landing]\label{thm:landing}
Suppose the zeros at time $t_0\ge0$ lie in $|\Ima z|\le\sqrt{w_0}$,
where $w_0>0$ and $t_0+w_0/2\le1/2$. With the parameters in
\eqref{eq:countparameters}, define
\begin{equation}\label{eq:landingfunction}
 \mathcal F_{n,B}(w)=\frac{w}{n+2}
 +\frac{B^2n}{4(n+2)^2}\log\left(1+\frac{2(n+2)w}{B^2}\right).
\end{equation}
Then
\begin{equation}\label{eq:landingbound}
 \Lambda\le t_0+\mathcal F_{n,B}(w_0)<t_0+\frac{w_0}{2}.
\end{equation}
More generally, if a continuous global floor $E_t(z)\ge g(w)\ge0$
holds at every simple maximizing zero throughout the relevant times,
the landing duration is at most
\begin{equation}\label{eq:generallanding}
 \int_0^{w_0}\frac{\dd u}{2+4u g(u)}.
\end{equation}
\end{theorem}
\begin{proof}
First take a compact positive-time interval. Uniform eventual reality
confines all its nonreal roots to a common compact rectangle; locally
uniform analyticity and the argument principle give a uniform finite
bound on their number there. The local Hermite expansion
\eqref{eq:hermite} separates the branches near every multiple root.
The same local argument applies at a complex multiple root: its
nonzero leading coefficient multiplies the heat evolution of
$(z-z_*)^m$, whose square-root rescaling has the same distinct Hermite
roots. The analytic implicit-function theorem after this rescaling
gives distinct branches analytic in the square-root time parameter.
Therefore multiple-root events in this compact region are locally
isolated, hence finite by compactness. Branches are analytic away
from these events, and analytic in a local square-root time parameter
near an event. Their derivatives have integrable singularities.

The maximal squared imaginary height $\mathcal W(t)$, with zero
included, is thus continuous and absolutely continuous on this
interval: locally it is the maximum of finitely many absolutely
continuous squared-height branches. At almost every time with
$\mathcal W>0$, a maximizing simple branch obeys
\[
 \mathcal W'\le-2-\frac{4n\mathcal W}{B^2+4\mathcal W}.
\]
All maximizers have the same field floor, so switches do not alter
the inequality. Direct differentiation gives
\begin{equation}\label{eq:landingderivative}
 \mathcal F_{n,B}'(w)
 =\frac{B^2+4w}{2B^2+4(n+2)w}
 =\frac1{2+4wn/(B^2+4w)}.
\end{equation}
Consequently $d\mathcal F_{n,B}(\mathcal W(t))/dt\le-1$ almost
everywhere while $\mathcal W>0$. Integrate across the finite
multiple-root events. If the maximal height has not vanished by
$t_0+\mathcal F_{n,B}(w_0)$, the integrated inequality contradicts
nonnegativity. Once it vanishes, upward persistence of real zeros
gives \eqref{eq:landingbound}. The same argument, with the integral
in \eqref{eq:generallanding} as the comparison function, proves the
general statement.

When $t_0=0$, start at $\varepsilon>0$ instead.
Strip contraction gives
$\mathcal W(\varepsilon)\le\max\{w_0-2\varepsilon,0\}$.
Apply the positive-time result and let $\varepsilon\downarrow0$.
This uses no attained nonreal maximum at time zero. The comparison
intervals stay within $t\le1/2$, since their duration is at most
the classical duration $w_0/2$.
\end{proof}

The strictly positive gain is explicit in the count parameters:
\begin{equation}\label{eq:landinggain}
 \frac w2-\mathcal F_{n,B}(w)
 =\int_0^w\frac{nu}{B^2+2(n+2)u}\dd u
 \ge\frac{nw^2}{2[B^2+2(n+2)w]}.
\end{equation}
It begins at $nw^2/(2B^2)$ as $w\downarrow0$. Applied to the
original unit strip this recovers the known qualitative conclusion
$\Lambda<1/2$ \cite{kikimlee}, without a new numerical endpoint. Making the
unprinted counting constant effective is a separate analytic task;
assigning an arbitrary number to the source's $O$-term is not a proof.

For a constant field floor $h>0$, the general duration is
\begin{equation}\label{eq:constantlanding}
 \mathcal T_h(w)=\frac{\log(1+2hw)}{4h}
 =\frac w2-\frac h2w^2+O(w^3).
\end{equation}
At a hypothetical ordinary collision from
Proposition~\ref{prop:double}, $t+\mathcal T_h(y(t)^2)\to T$.
Thus a bounded attraction field, including the count floor, is
compatible with a strictly positive limiting threshold.

% Consolidation of Heat Notes 8--10; insertion before the final analytic outlook.
% Prepared for Edward Baker with substantial LLM assistance, 9 October 2026.
% Model: GPT-6.1-sol (Codex); reasoning effort: ultra, verified from recorded configuration.

\section{The actual signed collision vector in shrinking time}
\label{sc:section}
The holomorphic interface of Section~\ref{sec:normalization} can be paid
uniformly in a shrinking-time regime. It then exposes a signed arithmetic
target, rather than a phase-independent tail estimate. We first record that
target and an actual signed edge reduction. Two phase-relaxation theorems
then delimit mechanisms that discard the genuine height correlations.
None of the results in this section supplies the missing local lower bound.

\subsection{Exact coordinates and complete disk payments}
Use
\begin{equation}\label{sc:scaling}
 1\le\kappa\le2,\quad 0<t\le1/20,\quad L=\kappa/t,\quad
 x=4\pi e^L,\quad R=L^{-1},\quad
 N=\left\lfloor\sqrt{e^L+t/16}\right\rfloor,
\end{equation}
and put
\begin{equation}\label{sc:exponents}
 \mathfrak a=\frac{\kappa(4-\kappa)}{16},\qquad
 \mathfrak b=\frac{\kappa(\kappa+4)}{16},\qquad
 \mathfrak d=\frac{\kappa(\kappa+12)}{16}.
\end{equation}
Thus $\mathfrak a+\mathfrak b=\kappa/2$,
$\mathfrak d-\mathfrak b=\kappa/2$, and
$\mathfrak d=\kappa-\mathfrak a$. All limits in this section are uniform
in $\kappa\in[1,2]$. Spatial derivatives hold time and the integer cutoff
fixed. When $L$ scales a derivative coordinate, it is the fixed center
value, not a quantity differentiated along \eqref{sc:scaling}.

\begin{proposition}[Complete shrinking-time error]
\label{sc:paiderror}
In \eqref{sc:scaling}, use $X=x$ and $\varepsilon=\tau=t$ in
\eqref{eq:diskconstants}--\eqref{eq:eta}, with the indicated natural
center cutoff $N$. Then
\begin{equation}\label{sc:diskpayment}
 \begin{aligned}
 \eta_N&\le5e^{-\mathfrak b/t},\qquad
 |Q_t(z)-F_{t,N}(z)|\le\eta_N\quad(|z-x|\le1/L),\\
 |Q_t^{(j)}(x)-F_{t,N}^{(j)}(x)|&\le j!L^j\eta_N
 \quad(j\ge0).
 \end{aligned}
\end{equation}
The estimates include reflection, conversion to the symmetric
normalizer, every possible cutoff change, and the spatial derivative errors.
\end{proposition}
\begin{proof}
We give the numerical ledger for the existing positive majorant. Here
$L\ge20$, $x>10^9$, $R\le1/20$, and $m_-\ge22000$.
The definitions in \eqref{eq:diskconstants} give
\[
 |\log q_\pm-L|<10^{-9},\quad
 |\alpha_- -L/2|<10^{-9},\quad
 A_*/L<.502,\quad D_*<1.01/x,\quad K_*<1.3.
\]
For example $|\log q_\pm-L|\le2R/x$, and
$K_*\le\exp\{RA_*(1+tD_*/2)/2\}$. There are at most two permitted
natural cutoffs, differing by at most one, and the center cutoff is one of
them. Every correction index $j$ satisfies
$|\log j-L/2|<10^{-4}$; substitution in \eqref{eq:termmajorants} gives
\[
 \log P_j\le-\mathfrak b/t+.251,\qquad
 2\sum_{\text{cutoff correction}}P_j\le2.6e^{-\mathfrak b/t}.
\]
The two positive corrections in $E_m^C$ together are less than $.001$;
the first is at most $2.5/(22000-.125)$, and the second follows from
the decrease of $L/e^L$ for $L\ge20$. Consequently
$E_m^C\le1.01e^{-\mathfrak b/t}$.

The many-term absolute sum used to pay the approximation error obeys
\[
 \sum_{n\le m}V_n\le5e^{\mathfrak a/t},\qquad
 1+m^{k_*}n^R<2.7,\qquad U_n\le.9/x.
\]
Indeed the decreasing-weight integral, with $n=e^{L/2-v}$, is at most
\[
 (1+10^{-8})e^{\mathfrak a/t}
 \int_0^{L/2}e^{-v/2+tv^2/4}\dd v
 \le(1+10^{-8})4e^{\mathfrak a/t}.
\]
Here $tv\le\kappa/2\le1$. Its first summand and possible edge interval
fit in the reserve to five because $\mathfrak a/t\ge15/4$.
The logarithmic weight slope is at most $-1/2+3\cdot10^{-6}$.
Also $n^R\le e^{1/2+5\cdot10^{-6}}$, $m^{k_*}\le e^{10^{-9}}$,
and $t^2\ell_n^2/16\le1/4+10^{-8}$; the numerator defining $U_n$
is at most $.87600001$. The resulting error bound is
\[
 E_m^{AB}\le\frac{5(2.7)(.9)}{4\pi}e^{-\mathfrak d/t}
 <e^{-\mathfrak d/t}\le.01e^{-\mathfrak b/t}.
\]
The last comparison uses $\kappa/(2t)\ge10$. Thus
\[
 \eta_N\le1.3(1.01+2.6+.01)e^{-\mathfrak b/t}<5e^{-\mathfrak b/t}.
\]
Proposition~\ref{prop:diskerror} now gives the full disk estimate and
Cauchy's formula gives every fixed-order center derivative. In particular
this proof applies arbitrarily near a natural-cutoff crossing; it never
differentiates that jumping cutoff.
\end{proof}

At the real center put $s=(1-ix)/2$, $\alpha=\alpha_r+i\alpha_i$,
and $\alpha'=U+iV$. Define the exact data
\begin{align}\label{sc:coefficients}
 w_n&=\exp\{t\log^2n/4-(1/2+t\alpha_r/2)\log n\},&
 \phi_n&=\theta_t(x)+(x-t\alpha_i)\log n/2,\\
 c&=\tfrac12(1+tU/2),&
 \Omega&=\tfrac12\Rea\{\alpha(1+t\alpha'/2)\},\nonumber\\
 \lambda_n&=\Omega-c\log n,&
 r_n&=4\lambda_n/L
      =\tfrac2L\Rea\{(\alpha-\log n)(1+t\alpha'/2)\},\nonumber\\
 c_n&=tV\log n/L,& z_n&=r_n-ic_n.\nonumber
\end{align}
These quantities retain the amplitude drift:
$(\log w_n)'=-tV\log n/4$ and $\phi_n'=-\lambda_n$.
In particular, with $h_n=(\cos\phi_n,r_n\sin\phi_n-c_n\cos\phi_n)$,
\begin{equation}\label{sc:vector}
 \mathcal V_N:=\left(F_{t,N}/2,2F_{t,N}'/L\right)
 =\sum_{n\le N}w_nh_n
 =\left(\Rea\sum_{n\le N}w_ne^{i\phi_n},
         \Ima\sum_{n\le N}w_nz_ne^{i\phi_n}\right).
\end{equation}
The term $c_n$ must not be removed when converting the derivative.
At a genuine collision, \eqref{eq:collisionnormalized} and
Proposition~\ref{sc:paiderror} imply
\begin{equation}\label{sc:collision}
 |\mathcal V_{N,1}|\le\eta_N/2,\qquad
 |\mathcal V_{N,2}|\le2\eta_N,\qquad
 |\mathcal V_N|^2\le17\eta_N^2/4
       \le\frac{425}{4}e^{-2\mathfrak b/t}.
\end{equation}
A uniform strict reverse of the squared bound for the actual phases would
exclude collisions in this regime. It has not been established.

The correlations required in such a bound can be written exactly. Put
$\Delta=\phi_n-\phi_m$ and $\Sigma=\phi_n+\phi_m$. Then
\begin{align}\label{sc:kernel}
 2h_n\cdot h_m={}&(1+r_nr_m+c_nc_m)\cos\Delta
 +(1-r_nr_m+c_nc_m)\cos\Sigma\nonumber\\
 &+(c_nr_m-c_mr_n)\sin\Delta
 -(c_mr_n+c_nr_m)\sin\Sigma.
\end{align}
Here $\Delta=(x-t\alpha_i)\log(n/m)/2$ and
$\Sigma=2\theta_t+(x-t\alpha_i)\log(nm)/2$.
Summing \eqref{sc:kernel} over all ordered pairs with weights $w_nw_m/2$
gives $|\mathcal V_N|^2$, including every diagonal and reflected
sum-frequency term.

For later estimates the exact real-axis formulas are
\begin{align}\label{sc:realaxis}
 \alpha_r&=L/2+\tfrac14\log(1+x^{-2})-(1+x^2)^{-1},&
 \alpha_i&=3x/(1+x^2)-\tfrac12\arctan x,\\
 U&=(7x^2-5)/(1+x^2)^2,&
 V&=x(x^2+5)/(1+x^2)^2.\nonumber
\end{align}
They imply, uniformly up to $N$,
\begin{equation}\label{sc:coordinateasymptotics}
 r_n=1-2\log n/L+O(t/(Lx)+1/(Lx^2)),\qquad c_n=O(t/x).
\end{equation}
\subsection{Absolute terminal-block budgets}

\begin{proposition}[Last-block absolute mass]\label{prop:absmass}
Take $x=4\pi e^{\kappa/t}$ and let
$N=\lfloor\sqrt{x/(4\pi)+t/16}\rfloor$. Uniformly for $\kappa$
in compact subsets of $(0,4)$, the real-axis coefficients in
\eqref{eq:p} satisfy
\begin{equation}\label{eq:lastblock}
 \sum_{N/2\le n\le N}|p_{t,n}((1-ix)/2)|
 =\left(2(1-2^{-1/2})+o(1)\right)
    \exp\left(\frac{\kappa(4-\kappa)}{16t}\right).
\end{equation}
At $\kappa=4$, a fixed logarithmic block
$Ne^{-M}\le n\le N$ instead has limiting absolute mass
$2(1-e^{-M/2})$ for each fixed $M>0$.
\end{proposition}
\begin{proof}
On the real axis a direct calculation from \eqref{eq:m0} gives
\[
 \Rea\alpha((1-ix)/2)
 =\tfrac12\log\frac{x}{4\pi}
      +\tfrac14\log(1+x^{-2})-\frac1{1+x^2}.
\]
Thus
$\log|p_n|=-(1/2+\kappa/4)\log n+(t/4)\log^2n+o(1)$,
with the error uniform on the indicated block. Write
$n=Ne^{-v}$, $0\le v\le\log2$, and use
$\log N=\kappa/(2t)+o(1)$. The weighted Riemann sum is
\[
 (1+o(1))e^{\kappa(4-\kappa)/(16t)}
 \int_0^{\log2}e^{-v/2}e^{tv^2/4}\dd v.
\]
Its integral tends to $2(1-2^{-1/2})$. The errors from $\alpha$
and the floor in $N$ are exponentially small in this scaling.
The same calculation on $[0,M]$ at $\kappa=4$ proves the final
assertion.
\end{proof}

This is a coefficient diagnostic, not a collision theorem for the
true heat function. It shows why treating the last block as a small
absolute perturbation of the leading term fails for $0<\kappa<4$.
A finite Dirichlet mollifier with constant coefficient one and
indices supported on a fixed finite set of primes cannot remove
this coefficient-wise absolute obstruction. In its convolution,
coefficients at integers coprime to those primes are unchanged;
their positive periodic density in the long block leaves the same
exponential scale in \eqref{eq:lastblock}. This observation excludes
that specific absolute method, not every mollifier or collision
argument. Genuine signed phase cancellation or the joint equations
must enter a successful proof in this range.

The last-block calculation in
Proposition~\ref{prop:absmass} is therefore not repaired by adding the
derivative coordinate. With $q=\log2$ its absolute scaled derivative
budget is
\begin{equation}\label{sc:derivativemass}
 \sum_{N/2\le n\le N}w_n\sqrt{r_n^2+c_n^2}
 =\left(\frac{2t}{\kappa}C_1+o(t)\right)e^{\mathfrak a/t},
 \quad C_1=\int_0^qv e^{-v/2}\dd v>0.
\end{equation}
To obtain this extension, use the same substitution as in that proof;
on its bounded logarithmic block $r_n=2tv/\kappa+o(t)$ and $c_n=o(t)$.
The extra polynomial factor $t$ does not overcome
$\mathfrak a\ge3/16$. This is an absolute-budget diagnostic, not a
lower bound on the actual signed derivative.

\subsection{An actual signed edge with explicit payment}
\begin{proposition}[The original value edge]
\label{sc:valueedge}
Let $H_v=\lfloor N^{3/4}\rfloor$ and $K_v=N-H_v$.
In \eqref{sc:scaling},
\begin{equation}\label{sc:edgetails}
 \left|\sum_{K_v<n\le N}w_ne^{i\phi_n}\right|\le40N^{-1/24},
 \qquad
 \left|\sum_{K_v<n\le N}w_nz_ne^{i\phi_n}\right|
     \le\frac{800}{L}N^{-7/24}.
\end{equation}
Consequently every genuine collision satisfies
\begin{equation}\label{sc:edgecollision}
 |F_{t,K_v}/2|\le\eta_N/2+40N^{-1/24},\qquad
 |2F_{t,K_v}'/L|\le2\eta_N+800N^{-7/24}/L.
\end{equation}
\end{proposition}
\begin{proof}
Put $T=(x-t\alpha_i)/2$ and $f(u)=T\log u/(2\pi)$.
The common phase does not affect a sum's modulus. Since
$N\ge22000$, $K_v\ge.9N$, and
$N^2-t/16\le e^L<(N+1)^2$, one has
\[
 1/N\le f'''(u)=T/(\pi u^3)\le4/N\quad(K_v\le u\le N).
\]
For the upper bound use $T/\pi\le2.01N^2$ and $2.01/.9^3<4$;
for the lower one use $T/\pi\ge N^2$.
The explicit theorem \cite[equation~(2)]{arias}, with derivative order
three, gives for a real smooth phase with $0<\lambda\le f'''\le\Lambda$
on an interval of length $Y$, $\lfloor Y\rfloor>3$,
\[
 \left|\sum e(f(n))\right|
 \le11\max\{Y^{3/4}(\Lambda/\lambda)^{1/4},
             Y(\Lambda^2/\lambda)^{1/6},Y^{1/4}\lambda^{-1/4}\}.
\]
Taking $\lambda=1/N$, $\Lambda=4/N$, every partial interval of this
edge is bounded by
\begin{equation}\label{sc:partialsum}
 11\max\{4^{1/4}N^{9/16},16^{1/6}N^{7/12},N^{7/16}\}
 \le18N^{7/12}.
\end{equation}
Intervals with at most three terms satisfy this estimate trivially.

Write $\delta=\log N-L/2$. The floor gives $|\delta|\le2/N$ and
\eqref{sc:realaxis} gives $|\alpha_r-L/2|\le x^{-2}$.
With $s_\kappa=1/2+\kappa/8$, the exact endpoint identity is
\[
 \log w_N+s_\kappa\log N
 =\frac t4\log N\,\delta-\frac t2(\alpha_r-L/2)\log N.
\]
Thus $w_N\le1.01N^{-s_\kappa}$. The logarithmic weight slope on
$[K_v,N]$ lies between $-.504$ and $-.499$; hence $w$ decreases and
$w_{K_v}\le2N^{-s_\kappa}$. Abel summation with
\eqref{sc:partialsum} bounds the value tail by
$36N^{7/12-s_\kappa}\le40N^{-1/24}$.

For the index interpolation of $z_n$, direct differentiation gives
\[
 |z_N|\le6/(LN),\qquad
 z'(u)=-\frac1u\left\{\frac2L(1+tU/2)+i\frac{tV}L\right\},
 \qquad |z'(u)|\le3/(Lu).
\]
Here the prime is an index derivative. For example the endpoint bound
uses $|U|\le8/x^2$, $|V|\le2/x$, $x\ge11N^2$, and
$|\alpha_r-\log N|\le2/N+x^{-2}$.
It follows that
\[
 \operatorname{TV}(z)\le4H_v/(LN),\qquad
 \sup|z|\le10H_v/(LN),
\]
\[
 |w_Nz_N|+\operatorname{TV}(wz)
 \le(6.06+28H_v)\frac{N^{-s_\kappa}}{LN}
 \le40\frac{H_vN^{-s_\kappa}}{LN}.
\]
Abel summation bounds the derivative sum by
$720N^{1/3-s_\kappa}/L\le800N^{-7/24}/L$.
This proves \eqref{sc:edgetails}; \eqref{sc:diskpayment} and
\eqref{sc:vector} give \eqref{sc:edgecollision} by exact signed
subtraction. The raw value and derivative tail payments are respectively
$80N^{-1/24}$ and $400N^{-7/24}$.
\end{proof}

The absolute mass of this same edge is
$(1+o(1))N^{1/4-\kappa/8}$, which grows for $\kappa<2$.
Thus Proposition~\ref{sc:valueedge} removes a genuinely growing signed
set of terms. Its retained core still contains $N-o(N)$ terms.
The displayed third-derivative estimate permits a value edge $N^\beta$
when $\beta<2/3+\kappa/8$, giving the uniform threshold $\beta<19/24$.
The larger derivative core and the derivative-order comparison below
refine this observable-specific statement.

\subsection{A long-block exponent-pair budget}
\label{sc:longblock}
For a logarithmic-phase exponent pair $(p,q)$, its partial-interval bound
on a dyadic block near $M=e^{\vartheta L}$ is
$O(T^pM^{q-p})$; the phase meets the logarithmic case of the definition
in \cite[Sections~1.1--1.2]{trudgianyang}.
Abel summation with the exact decreasing weight gives
\begin{equation}\label{sc:pairbudget}
 |\text{weighted block}|
 \le\exp\{L(E_{p,q}(\vartheta,\kappa)+o(1))\},\qquad
 E_{p,q}=p+\vartheta(q-p-1/2-\kappa/4)+\kappa\vartheta^2/4.
\end{equation}
At the endpoint block $\vartheta=1/2$, a strictly decaying bound of
this form requires
\begin{equation}\label{sc:pairthreshold}
 p+q<1/2+\kappa/8.
\end{equation}
The Bourgain pair in the cited survey is
$(13/84+\epsilon,55/84+\epsilon)$; its endpoint exponent is
$13/84-\kappa/16+\epsilon\ge5/168+\epsilon>0$.
More generally, the particular convex hull $\mathscr H$ defined in
\cite[Section~1.4]{trudgianyang} has
\begin{equation}\label{sc:surveyedminimum}
 \min_{(p,q)\in\mathscr H}(p+q)=17/21.
\end{equation}
Here is the elementary minimum check, so the scope of the diagnostic is
explicit. The first five nonnegative-index vertices are
\[
 (13/84,55/84),\quad(4742/38463,35731/51284),\quad
 (18/199,593/796),\quad(2779/38033,58699/76066),\quad
 (715/10238,7955/10238).
\]
Their sums are at least $17/21$ by direct rational comparison, with
equality at the first. The next four are the $A$-transforms of the last
four displayed vertices, where
$A(p,q)=(p/(2p+2),(p+q+1)/(2p+2))$.
Each has $q\ge2/3$, so its transformed sum is at least $5/6>17/21$.
The negative-index vertices are $B$-reflections, which preserve $p+q$.
For the remaining family indexed by integers $m\ge5$, its deficit from
one is
\[
 \frac{3m^2-7m+2}{m(m-1)^2(m+2)}
 \le\frac3{(m-1)(m+2)}\le\frac3{28}<\frac4{21}.
\]
The limiting endpoints have sum one, and convexity preserves this linear
minimum. This proves \eqref{sc:surveyedminimum} for that specified hull,
which misses \eqref{sc:pairthreshold} throughout the present range.
It is not a claim about every later exponent-pair result or every signed
method. A positive exponent in an upper bound is a deficit of the bound,
not a lower bound asserting growth of the actual signed block.

\subsection{A genuine adjacent packet defeats diagonal coercivity}
\begin{proposition}\label{sc:packet}
Let $M\to\infty$ be integral, $x_M=4\pi M^2$, and
$t_M=\kappa/(2\log M)$, $1\le\kappa\le2$.
The natural cutoff is $N=M$. For its two genuine terms put
\[
 D_M=|w_Mh_M|^2+|w_{M-1}h_{M-1}|^2,\qquad
 J_M=|w_Mh_M+w_{M-1}h_{M-1}|^2.
\]
Then uniformly in $\kappa$,
\begin{equation}\label{sc:packetenergy}
 D_M=(2\cos^2(\pi/8)+o(1))w_M^2,\qquad
 J_M=O(w_M^2/M^2),\qquad J_M/D_M\longrightarrow0.
\end{equation}
The same ratio tends to zero after summing over any fixed finite collection
of bounded spatial translations, with time and the two indices fixed.
\end{proposition}
\begin{proof}
The exact branch of \eqref{eq:symmetricnormalization} gives
\[
 \theta_t=-\pi+\tfrac14\arctan x+\tfrac x4(1-L)
 -\tfrac x8\log(1+x^{-2})+\tfrac t2\alpha_r\alpha_i.
\]
At the specified points $L=2\log M$, so
\[
 \phi_M=\pi M^2-7\pi/8+O(M^{-2}),\qquad
 \phi_M-\phi_{M-1}=2\pi M+\pi+O(M^{-1}).
\]
For the second expansion use
$-\log(1-1/M)=M^{-1}+(2M^2)^{-1}+O(M^{-3})$ and
$T=2\pi M^2+O(t)$.
The two cosines have opposite signs up to $O(M^{-1})$, and each squared
cosine tends to $\cos^2(\pi/8)$. The exact weights satisfy
\[
 w_M=M^{-(4+\kappa)/8}(1+o(1)),\qquad
 w_{M-1}/w_M=1+O(M^{-1}).
\]
Equation~\eqref{sc:coordinateasymptotics} gives
$r_M=O(t/(LM^2)+1/(LM^4))$, $r_{M-1}=O(1/(ML))$,
and $c_n=O(t/M^2)$. The packet's value coordinate is therefore
$O(w_M/M)$ and its scaled derivative coordinate is $O(w_M/(ML))$.
This proves \eqref{sc:packetenergy}; its cross term equals
$-D_M+O(w_M^2/M^2)$.
For a bounded translation $y$, exact spatial differentiation gives
$\phi_M'=O(M^{-2})$ and $\phi_{M-1}'=O(M^{-1})$ in the corresponding
neighborhood; the weight and derivative-coordinate bounds remain uniform.
For finitely many such probes, the total packet energy is still
$O(w_M^2/M^2)$, while the total diagonal has the stated positive leading
coefficient multiplied by the number of probes.
\end{proof}
This disproves a positive constant packetwise lower bound by separate
diagonal energies, even with genuine phases and fixed bounded probes.
The packet is itself on the small cutoff-coefficient scale. The statement
does not concern cancellation of the complete sum or exclude larger-range
probes that use correlations between different packets.

\subsection{Independent coefficient phases can cancel the exact joint vector}
\begin{proposition}\label{sc:independent}
For all sufficiently small $t>0$, uniformly in $\kappa\in[1,2]$, there
exist phases $\psi_n$, with $\psi_1=\theta_t(x)$, such that
\begin{equation}\label{sc:independentzero}
 \sum_{n\le N}w_n(\cos\psi_n,r_n\sin\psi_n-c_n\cos\psi_n)=0.
\end{equation}
Every weight and local derivative coefficient is the exact one in
\eqref{sc:coefficients}.
\end{proposition}
\begin{proof}
Uniformly as $t\to0$, fixed-index weights tend to $n^{-\sigma}$ with
$\sigma=1/2+\kappa/4\in[3/4,1]$. The exact weight slope is negative up
to $N$. Equations~\eqref{sc:realaxis}--\eqref{sc:coordinateasymptotics}
therefore give, for sufficiently small $t$,
\[
 \max|r_n|\le1.01,\quad w_{257}\le.02,\quad |r_1-1|\le.005,
 \quad\max_{n\le256}\sqrt{c_n^2+(r_n-1)^2}\le.01,
\]
\[
 4.5\le S:=\sum_{n=2}^{256}w_n\le16,\qquad w_2\le.61.
\]
For the head-sum margins use $H_{256}-1>4.5$ and
$\int_1^{256}u^{-3/4}\dd u=12$.
For $n>256$ choose $\psi_n=\pm\pi/2$. Greedy signing of the real
numbers $w_nr_n$ leaves a derivative residual $d$ with
$|d|\le.0202<.025$; both value and amplitude-drift coordinates vanish.
The target including the fixed leading term is
\[
 T_0=(\cos\theta_t,r_1\sin\theta_t+d),\qquad .97\le|T_0|\le1.03.
\]
Partition the remaining head greedily into groups of masses $A,B$, putting
each next weight in the lighter group. Then $A+B=S$, $|A-B|\le.61$,
and $A,B\ge1.945>|T_0|$. Common phases in the groups give vectors
$L_Ae(u),L_Be(v)$, $e(u)=(\cos u,\sin u)$, with
\[
 \|L_A-AI\|\le.01A,\qquad \|L_B-BI\|\le.01B.
\]
Here the exact first row is $(A,0)$ and the second row is
$(-\sum_Aw_nc_n,\sum_Aw_nr_n)$, and similarly for $B$.
Put $r=|T_0|$. Since $r-|A-B|\ge.36>.01(A+B)$, the quantity
$|L_B^{-1}(T_0+L_Ae(u))|$ is above one at $e(u)=T_0/r$ and below
one at $e(u)=-T_0/r$. The singular values of $L_B$ lie between
$B(1-.01)$ and $B(1+.01)$, which verifies both strict comparisons.
Continuity supplies a crossing; choose
$e(v)=-L_B^{-1}(T_0+L_Ae(u))$ there. This cancels the entire vector.
\end{proof}
Multiplying the first analytic term by a constant phase and its reflection
by the opposite phase realizes this construction with real symmetry and
unchanged local derivative data. These phases generally violate
$\phi_n=\theta_t+T\log n$. Proposition~\ref{sc:independent} is a
theorem about a relaxation, not a collision of $F_{t,N}$ or $H_t$.

\subsection{Even complete multiplicativity does not force a positive vector}
\begin{theorem}[Complete multiplicative-phase obstruction]
\label{sc:character}
There exists $t_0>0$ such that at every parameter point
$0<t<t_0$, $\kappa\in[1,2]$, a completely multiplicative function
$\chi:\mathbb N\to\mathbb C$ with $\chi(1)=1$ and $|\chi(n)|=1$
satisfies
\begin{equation}\label{sc:characterzero}
 \mathcal V_\chi:=\sum_{n\le N}w_n
  (\cos\psi_n,r_n\sin\psi_n-c_n\cos\psi_n)=0,
 \qquad\psi_n=\phi_n+\arg\chi(n).
\end{equation}
All terms are retained and the leading phase is fixed. No explicit value
of $t_0$, common twist for a neighborhood, or genuine-height realization
is asserted.
\end{theorem}
\begin{proof}
We select a bounded complete residual, then cancel it using isolated
primes. The coefficient formulas give, uniformly for small $t$,
\begin{equation}\label{sc:charactersquarebound}
 \max_{n\le N}|r_n|\le1.01,\quad\max|c_n|\le.01,
 \qquad w_n^2\le n^{-6/5}.
\end{equation}
For the last bound, the coefficient of $\log n$ in $\log w_n^2$ is at
most $-(1+t\alpha_r)+(t/2)\log N=-1-\kappa/4+o(1)$.
The logarithmic slope of $w$ is at most $-.49$ up to $N$.

Select $\mathcal P=\{p\text{ prime}:N/2<p\le3N/4\}$.
Since $2p>N$, each selected prime divides only its own term among the
retained integers. Its twist phase can be changed independently while
preserving complete multiplicativity of the whole twist.
Assign independent Haar unit phases to every unselected prime at most
$N$, and extend them multiplicatively to the retained integers outside
$\mathcal P$. Their complete vector sum, including $n=1$, is $T_0$.
Unique factorization gives
\[
 \mathbb E\chi(n)\overline{\chi(m)}=\mathbf1_{n=m},\qquad
 \mathbb E\chi(n)\chi(m)=0
\]
for these indices, with the second identity excepting $n=m=1$.
Every nonzero prime-exponent vector is annihilated by Haar integration.
It follows that every cross term is paid and
\[
 \mathbb E|T_0|^2
 =|(\cos\theta_t,r_1\sin\theta_t)|^2
 +\frac12\sum_{\substack{2\le n\le N\\n\notin\mathcal P}}
    w_n^2(1+r_n^2+c_n^2)
 \le1.0201+1.0101\int_1^\infty u^{-6/5}\dd u<7.
\]
Fix a realization with $|T_0|\le\sqrt7$. No selected-prime phase has
entered that sum.

For $p=Nz$, $1/2<z\le3/4$, direct substitution gives
$w_p=e^{-\mathfrak b/t}z^{-1/2}(1+o(1))$ uniformly.
Import only the classical prime number theorem
$\pi(y)\sim y/\log y$, as recorded in
\href{https://dlmf.nist.gov/27.12\#E4}{NIST DLMF, equation 27.12.4}.
Partial summation on this fixed relative interval gives
\begin{equation}\label{sc:primereservoir}
 S:=\sum_{p\in\mathcal P}w_p
 =\left(\frac{2t}{\kappa}(\sqrt3-\sqrt2)+o(t)\right)e^{\mathfrak a/t},
 \qquad tS\longrightarrow\infty.
\end{equation}
Uniformity follows because the smallest permitted $N$ tends to infinity.
Furthermore
$r_p/t=2\log(N/p)/\kappa+o(1)$, so
\[
 t/4\le r_p\le1.5t,\qquad c_p=O(t/x).
\]
The three sequences $w_p,w_pr_p,w_pc_p$, in increasing prime order,
are positive and decreasing. The second assertion uses the exact affine
decrease of $r_p$ in $\log p$. For the third, $V>0$ and its
logarithmic derivative has sign
$\{d\log w/d\log p\}\log p+1<0$ for small $t$.

Assign the selected primes cyclically to three groups. For a decreasing
nonnegative sequence $a_1,\ldots,a_m$, their group sums differ by at
most $a_1$: extend by zeros and telescope
\[
 \sum_{j\ge0}(a_{3j+1}-a_{3j+3})
 \le\sum_{j\ge0}(a_{3j+1}-a_{3j+4})=a_1.
\]
Apply this to each of the three sequences. Giving group $j$ a common
total phase $\beta_j$ yields the exact vector $L_je(\beta_j)$, with
\[
 L_j=\begin{pmatrix}\sum_jw_p&0\\-\sum_jw_pc_p&\sum_jw_pr_p\end{pmatrix},
 \quad B=\frac13\sum_jL_j=\frac13\begin{pmatrix}S&0\\-C&D\end{pmatrix},
\]
where $C=\sum_{\mathcal P}w_pc_p$, $D=\sum_{\mathcal P}w_pr_p$.
Since $D/S\in[t/4,1.5t]$ and $C/S=O(t/x)$,
\[
 B^{-1}=\frac3S\begin{pmatrix}1&0\\C/D&S/D\end{pmatrix},
 \quad\|B^{-1}\|=O((tS)^{-1}),\quad
 \max_j\|L_j-B\|=O(e^{-\mathfrak b/t}).
\]
Consequently, with Euclidean and induced matrix norms,
\begin{equation}\label{sc:contractionsmallness}
 \epsilon:=\max_j\|B^{-1}(L_j-B)\|
 =O(e^{-\mathfrak b/t}/(tS))\longrightarrow0,
 \qquad\tau:=|B^{-1}T_0|=O((tS)^{-1})\longrightarrow0.
\end{equation}

For completeness an explicit uniform contraction solves the two real
coordinates exactly. Put $q_1=0$, $q_2=2\pi/3$, $q_3=4\pi/3$,
hold $\beta_3=q_3$, and set $\beta_1=z_1$, $\beta_2=q_2+z_2$.
The equation is
\[
 h(z)+E(z)+B^{-1}T_0=0,
 \quad h(z)=e(z_1)+e(q_2+z_2)+e(q_3),
 \quad E(z)=\sum_jB^{-1}(L_j-B)e(\beta_j).
\]
At zero,
\[
 J=Dh(0)=\begin{pmatrix}0&-\sqrt3/2\\1&-1/2\end{pmatrix},
 \qquad\|J^{-1}\|=\sqrt2.
\]
Writing $h(z)=Jz+R(z)$ gives
$|R(z)|\le|z|^2/2$, $\|DR(z)\|\le|z|$,
$|E(z)|\le3\epsilon$, and $\|DE(z)\|\le\sqrt2\epsilon$.
For small enough $t$, take $\epsilon\le1/100$, $\tau\le1/50$.
On $|z|\le1/10$, the map
$-J^{-1}\{R(z)+E(z)+B^{-1}T_0\}$ has norm at most
\[
 \frac32\left(\frac1{200}+\frac3{100}+\frac1{50}\right)
 =\frac{33}{400}<\frac1{10},
\]
and Lipschitz constant at most
$\frac32\frac1{10}+2/100=17/100<1$.
Banach's theorem supplies the required exact solution.
For a selected prime in group $j$, set
$\chi(p)=e^{i(\beta_j-\phi_p)}$; keep the fixed unselected prime phases,
assign unused prime phases arbitrarily, and extend completely
multiplicatively. Isolation ensures this final assignment alters only
the selected-prime terms among all $n\le N$. It proves
\eqref{sc:characterzero} with no omitted composite or small-prime term.
\end{proof}

The heat weight has not been made multiplicative: its exact cross factor
$e^{(t/2)\log m\log n}$ remains when indices are multiplied.
Only $\chi$ is completely multiplicative. Multiplying first analytic
terms by $\chi(n)$ and their reflections by $\overline{\chi(n)}$
preserves real symmetry and, with the twist held fixed, every local
amplitude derivative and phase frequency. No approximation theorem for
a twisted heat function is asserted. The genuine case is $\chi\equiv1$;
the constructed twists generally depart from its logarithmic height orbit.
Approximating finitely many prime phases by a different height does not
transfer the theorem: that movement also changes the cutoff, weights,
normalizer, carrier, and derivative coefficients. These joint changes
have not been paid by the twist construction.

\subsection{The local input that remains}
\label{sc:remaining}
Writing $D=\sum_nw_n^2|h_n|^2$, an equivalent sufficient signed target
for the full vector is
\begin{equation}\label{sc:signedtarget}
 2\sum_{n<m}w_nw_mh_n\cdot h_m>-D+17\eta_N^2/4.
\end{equation}
For the original edge core one may instead prove a strict reverse of
either bound in \eqref{sc:edgecollision}. Both targets must retain the
actual shared-height phases and every stated error.
Propositions~\ref{sc:packet} and \ref{sc:independent}, and
Theorem~\ref{sc:character}, respectively rule out packetwise diagonal
coercivity and positive joint lower bounds valid for all independent or
all completely multiplicative unit phases. They do not rule out a
height-conditioned signed estimate for the genuine sum.

The earlier complete value-coordinate probe used the sufficient range
$H_p=CLe^{2\mathfrak a/t}$. The weighted derivative probe in
Theorem~\ref{lj:derivativeaverage} below supplies a shorter sufficient
range and pays the genuine physical movement and heat remainder.
We therefore use that stronger averaged result instead of duplicating
the earlier mean-square proof. An average alone still permits isolated
double zeros. The following sections quantify its remaining local
transfer deficit and add an all-real-threshold jet constraint. The
regime \eqref{sc:scaling} also leaves other scales and low heights to
control; no global collision exclusion or RH conclusion follows here.


% Integration insert for Heat Notes 11--12, 9 October 2026.
% Prepared for Edward Baker with substantial LLM assistance.
% Model: GPT-6.1-sol (Codex); reasoning effort: ultra, verified from recorded configuration.
% Internal same-model audits are not independent validation.
% This fragment has no preamble. It follows the signed shrinking-time section.

\section{Asymmetric signed cores and complete derivative averages}
\label{lj:cores}

We retain the scaling \eqref{sc:scaling}, the symmetric normalizer
\eqref{eq:symmetricnormalization}, and the exact genuine coefficients from
the preceding section. In particular,
\[
 L=\kappa/t,\qquad x=4\pi e^L,\qquad
 N=\left\lfloor\sqrt{e^L+t/16}\right\rfloor,\qquad
 1\le\kappa\le2,\quad 0<t\le1/20,
\]
\[
 \mathfrak a=\frac{\kappa(4-\kappa)}{16},\qquad
 \mathfrak b=\frac{\kappa(\kappa+4)}{16},\qquad
 s_\kappa=\frac12+\frac{\kappa}{8}.
\]
Unless stated otherwise, limits and implicit constants are uniform in this
range. Every spatial derivative holds time and its integer cutoff fixed.
When \(L\) scales a derivative coordinate at a translated point, it is the
fixed value at the original center.

Write \(s=(1-ix)/2\), \(\alpha(s)=\alpha_r+i\alpha_i\), and
\(\alpha'(s)=U+iV\). The exact summands \(q_n=w_ne^{i\phi_n}\) satisfy
\begin{align}
 w_n&=\exp\left\{\frac t4\log^2n-
             \left(\frac12+\frac{t\alpha_r}{2}\right)\log n\right\},
 &\phi_n&=\theta_t(x)+\frac{x-t\alpha_i}{2}\log n,
 \label{lj:exactdata}\\
 c&=\frac12(1+tU/2),
 &\Omega&=\frac12\Rea\{\alpha(1+t\alpha'/2)\},\nonumber\\
 \lambda_n&=\Omega-c\log n,\qquad r_n=4\lambda_n/L,
 &c_n&=tV\log n/L,\qquad z_n=r_n-ic_n.\nonumber
\end{align}
Thus the real-axis identities, including amplitude differentiation, are
\begin{equation}\label{lj:exactcoordinates}
 \frac{F_{t,N}}2=\Rea\sum_{n\le N}q_n,\qquad
 \frac{2F_{t,N}'}L=\Ima\sum_{n\le N}q_nz_n.
\end{equation}
The complete radius-\(1/L\) holomorphic payment from the preceding
section gives
\begin{equation}\label{lj:cauchypayment}
 \eta_N\le5e^{-\mathfrak b/t},\qquad
 |Q_t^{(j)}(x)-F_{t,N}^{(j)}(x)|\le j!L^j\eta_N.
\end{equation}
This payment already includes conversion to the symmetric normalizer,
reflection, and every possible natural-cutoff change.

\subsection{A larger signed edge for the derivative coordinate}

Put
\[
 H_{e,v}=\lfloor N^{3/4}\rfloor,\quad K_v=N-H_{e,v},
 \qquad H_{e,d}=\lfloor N^{7/8}\rfloor,\quad K_d=N-H_{e,d}.
\]
The signed value estimate in Proposition~\ref{sc:valueedge} retains \(K_v\).
The derivative coordinate permits the wider removal \(N-K_d\).

\begin{proposition}[Asymmetric derivative core]\label{lj:derivativecore}
In the stated range,
\begin{equation}\label{lj:derivativetail}
 \left|\sum_{K_d<n\le N}w_nz_ne^{i\phi_n}\right|
       \le\frac{800}{L}N^{-1/24},\qquad
 |F_{t,N}'-F_{t,K_d}'|\le400N^{-1/24}.
\end{equation}
Consequently every genuine collision satisfies
\begin{equation}\label{lj:mixedcollision}
 \left|\frac{F_{t,K_v}}2\right|
       \le\eta_N/2+40N^{-1/24},\qquad
 \left|\frac{2F_{t,K_d}'}L\right|
       \le2\eta_N+\frac{800}{L}N^{-1/24}.
\end{equation}
The two quantities in \eqref{lj:mixedcollision} have different fixed
cutoffs; the second is not the derivative of the first.
\end{proposition}

\begin{proof}
The scaling implies \(N\ge22000>(10/3)^8\), hence \(K_d\ge.7N\).
Set
\[
 T=\frac{x-t\alpha_i}{2},\qquad f(u)=\frac{T\log u}{2\pi}.
\]
On \(K_d\le u\le N\), the exact real logarithmic phase obeys
\[
 \frac1N\le f'''(u)=\frac{T}{\pi u^3}\le\frac6N.
\]
Indeed \(N^2\le T/\pi\le2.01N^2\) and \(2.01/.7^3<6\).
The explicit third-derivative theorem
\cite[equation (2)]{arias}, with lower and upper derivative bounds
\(1/N\) and \(6/N\), bounds every partial interval of length at most
\(N^{7/8}\) by
\begin{equation}\label{lj:thirdpartials}
 \left|\sum e(f(n))\right|
 \le11\max\{6^{1/4}N^{21/32},
           36^{1/6}N^{17/24},N^{15/32}\}
 \le20N^{17/24}.
\end{equation}
The theorem requires the integer part of the interval length to exceed
three; the remaining partial intervals are bounded trivially.
The constant reserves are
\(6\cdot11^4<20^4\), \(36\cdot11^6<20^6\), and \(11<20\).

For the continuous index interpolation, the exact endpoint estimates are
\[
 w_N\le1.01N^{-s_\kappa},\qquad
 |z_N|\le\frac6{LN},\qquad
 \left|\frac{dz(u)}{du}\right|\le\frac3{Lu}.
\]
The logarithmic slope
\[
 \frac{d\log w(u)}{d\log u}
       =\frac t2\log u-\frac12-\frac{t\alpha_r}{2}
\]
lies between \(-.511\) and \(-.499\) on \([.7N,N]\).
Here one uses \(|\log N-L/2|\le2/N\),
\(|\alpha_r-L/2|\le x^{-2}\), and \(|\log .7|<.4\).
In particular \(w\) decreases and
\(w_{K_d}\le(10/7)w_N\le2N^{-s_\kappa}\). Its index variation and the
variation of \(z\) give
\[
 \operatorname{TV}(z)<\frac{5H_{e,d}}{LN},\qquad
 \sup|z|\le\frac{11H_{e,d}}{LN},
\]
\begin{equation}\label{lj:productvariation}
 |w_Nz_N|+\operatorname{TV}(wz)
 \le(6.06+32H_{e,d})\frac{N^{-s_\kappa}}{LN}
 \le40\frac{H_{e,d}N^{-s_\kappa}}{LN}.
\end{equation}
For the product estimate, use
\(\operatorname{TV}(wz)\le
 w_{K_d}\operatorname{TV}(z)+\sup|z|\operatorname{TV}(w)\)
and \(\operatorname{TV}(w)\le2N^{-s_\kappa}\).
Abel summation with \eqref{lj:thirdpartials} therefore gives
\[
 \left|\sum_{K_d<n\le N}q_nz_n\right|
 \le\frac{800}{L}N^{17/24+7/8-1-s_\kappa}
 =\frac{800}{L}N^{7/12-s_\kappa}
 \le\frac{800}{L}N^{-1/24}.
\]
Identity \eqref{lj:exactcoordinates} proves the raw derivative bound.
At a collision \(Q_t=Q_t'=0\), combine
\eqref{lj:cauchypayment} with this exact signed subtraction and the
value edge Proposition~\ref{sc:valueedge}. This proves
\eqref{lj:mixedcollision}.
\end{proof}

\subsection{Exact raw jets and a diagnostic for fixed derivative orders}

Raw higher derivatives must be estimated directly; differentiating an
edge length or a parameter-dependent cutoff would change the question.
Define the exact logarithmic derivative
\begin{equation}\label{lj:rawmultiplier}
 v_n=\frac{q_n'}{q_n}
       =-\frac{tV\log n}{4}-i\lambda_n.
\end{equation}
If written as \(-tV\log n/4-iLr_n/4\), the last term still means
the unscaled expression
\(\frac12\Rea\{(\alpha-\log n)(1+t\alpha'/2)\}\);
it is differentiated with \(t,n\) fixed.

\begin{proposition}[Fixed-time raw derivative tails]\label{lj:rawtails}
For every fixed \(j\ge1\),
\begin{align}
 |F_{t,N}^{(j)}-F_{t,K_v}^{(j)}|
    &=O_j(N^{7/12-s_\kappa-j/4}),\label{lj:shortjets}\\
 |F_{t,N}^{(j)}-F_{t,K_d}^{(j)}|
    &=O_j(N^{17/24-s_\kappa-j/8}).\label{lj:deepjets}
\end{align}
For \(j=2,3,4\), their uniform exponents are respectively
\((-13/24,-19/24,-25/24)\) and
\((-1/6,-7/24,-5/12)\).
The constants in these higher-jet bounds are uniform constants; the
explicit first-derivative certificate is
\eqref{lj:derivativetail}.
\end{proposition}

\begin{proof}
On an endpoint edge of length \(H=\lfloor N^\beta\rfloor\), with fixed
\(0<\beta<1\), put \(h=H/N\). The real-axis formulas for \(\alpha\)
and the endpoint frequency estimate give
\begin{equation}\label{lj:variationjets}
 \sup|v_n|+\operatorname{TV}_n(v_n)=O(h),\qquad
 \sup|\partial_x^kv_n|=O_k(N^{-2k}),\qquad
 \operatorname{TV}_n(\partial_x^kv_n)=O_k(hN^{-2k})
 \quad(k\ge1).
\end{equation}
Here \(\operatorname{TV}_n\) varies the continuous index at fixed
\(t,x\). To check the spatial estimates, use
\[
 \partial_x^k\alpha_r=O_k(x^{-k}),\quad
 \partial_x^k\alpha_i=O_k(x^{-k-1}),\quad
 \partial_x^kU=O_k(x^{-k-2}),\quad
 \partial_x^kV=O_k(x^{-k-1})\quad(k\ge1),
\]
with \(t\log n=O(1)\) and \(x\asymp N^2\).
Each \(\partial_x^kv_n\) is affine in \(\log n\); its varying
coefficient is \(O_k(x^{-k})\). The logarithmic length of the edge
is \(O(h)\), proving its total-variation estimate. For \(k=0\),
the index slope of \(\lambda_n\) is \(O(1/n)\), and
\(|v_N|=O(1/N)\).

Write \(q_n^{(j)}=q_nP_j\). The complete derivative polynomials through
the order used below are
\begin{align}
 P_1&=v,\qquad P_2=v^2+v',\qquad
 P_3=v^3+3vv'+v'',\nonumber\\
 P_4&=v^4+6v^2v'+3(v')^2+4vv''+v'''.\label{lj:bell}
\end{align}
Every monomial has differential weight \(j\), assigning weight \(k+1\)
to \(v^{(k)}\). Since \(N^{-2k}=O(h^{k+1})\) for \(k\ge1\),
product variation in \eqref{lj:variationjets} gives
\[
 \sup|P_j|+\operatorname{TV}_n(P_j)=O_j(h^j),\qquad
 |w_NP_j(N)|+\operatorname{TV}_n(wP_j)
             =O_j(N^{-s_\kappa}h^j).
\]
All amplitude and frequency derivatives in \eqref{lj:bell} are included.
For \(H=N^{7/8}+O(1)\), use \eqref{lj:thirdpartials}; for
\(H=N^{3/4}+O(1)\), the same source theorem bounds every partial sum
by \(O(N^{7/12})\). Indeed its three powers are
\(9/16,7/12,7/16\), and short intervals are again trivial.
Abel summation proves \eqref{lj:shortjets}--\eqref{lj:deepjets}.
\end{proof}

The calculation admits a precise limitation within the imported
fixed-order derivative theorem. For a fixed \(d\ge3\), the phase satisfies
\[
 f^{(d)}(u)=(-1)^{d-1}\frac{(d-1)!T}{2\pi u^d}
                   \asymp_d N^{2-d}.
\]
For even \(d\), conjugate the sum and apply the theorem to \(-f\).
Put \(A_d^{\rm pow}=2^{1-d}\) and
\(B_d^{\rm pow}=(d-2)/(2^d-2)\). These denote exponents, not the
constants named \(A_d,B_d\) in \cite{arias}. Its three terms, multiplied
by the interval length, give the partial-sum exponent
\begin{equation}\label{lj:orderdiagnostic}
 v_d(\beta)=\max\{\beta(1-A_d^{\rm pow}),\
                 \beta-B_d^{\rm pow},\
                 \beta(1-dA_d^{\rm pow})+(d-2)A_d^{\rm pow}\}.
\end{equation}
All interval-length powers are nonnegative. The same upper bound thus
holds for every partial interval, including the trivial short ones.
The raw-\(j\) tail bound supplied by this calculation is
\begin{equation}\label{lj:generalrawtail}
 O_{d,j}\bigl(N^{v_d(\beta)-s_\kappa+j(\beta-1)}\bigr).
\end{equation}

For \(d=3\), the middle term determines its strict decay threshold:
\begin{equation}\label{lj:rawthreshold}
 \beta<\frac{j+s_\kappa+1/6}{j+1}
       =1-\frac{5/6-s_\kappa}{j+1};
 \qquad \beta<1-\frac{5}{24(j+1)}
       \quad\hbox{uniformly}.
\end{equation}
This includes \(j=0\), for which the uniform value threshold is
\(\beta<19/24\). For completeness, the cross-differences between
the other two thresholds and the middle one have positive numerators
\((2j+6s_\kappa-3)/24\) and
\(j/3+3s_\kappa/4-7/24\).
For \(d\ge4\), the first term alone caps the threshold by
\((j+s_\kappa)/(j+7/8)\). It is smaller than
\eqref{lj:rawthreshold}, since the opposite cross-difference is
\((2j-6s_\kappa+7)/48>0\). Thus order three is optimal for each
fixed raw derivative order within this derivative-bound and Abel
calculation.

The value estimate also has a macroscopic-block deficit:
\(v_3(1)=5/6\), whereas \(v_d(1)\ge7/8\) for \(d\ge4\), and
\(5/6-s_\kappa\ge1/12>0\). Splitting a fixed-proportion endpoint
block into \(N^{1-\beta}\) pieces and adding these upper bounds
retains the exponent \(1-B_3^{\rm pow}=5/6\); for \(d\ge4\) it
retains \(1-\beta A_d^{\rm pow}\ge7/8\).
This is a deficit of the displayed upper bounds and triangle
aggregation. It is not a lower bound on a genuine sum or an obstruction
to transformations retaining correlations between pieces.

\subsection{Complete derivative mean square without a logarithmic loss}

The exact derivative weights near the cutoff improve a sufficient
averaging range by a factor \(t^2\). This is an averaged statement,
including every coefficient of the original approximant.

\begin{theorem}[Complete genuine derivative average]
\label{lj:derivativeaverage}
There are absolute \(D,t_0>0\) such that, for \(0<t<t_0\), put
\[
 H_d=Dt^2e^{2\mathfrak a/t},\qquad
 G_F(y)=2F_{t,N}'(x+y)/L,\qquad G_Q(y)=2Q_t'(x+y)/L.
\]
Then
\begin{equation}\label{lj:averagelower}
 \frac1{H_d}\int_0^{H_d}G_F(y)^2\dd y\ge\frac14,\qquad
 \frac1{H_d}\int_0^{H_d}G_Q(y)^2\dd y\ge\frac18.
\end{equation}
The corresponding complete joint energies, obtained by adding
\((F_{t,N}(x+y)/2)^2\) or \((Q_t(x+y)/2)^2\), inherit these lower
bounds. The denominator \(L\) is fixed at the original center.
\end{theorem}

\begin{proof}
For distinct real frequencies \(\omega_j\), let
\(\delta_j=\min_{k\ne j}|\omega_j-\omega_k|\).
The weighted Montgomery--Vaughan inequality
\cite[Theorem 2, equation (1.7), and Corollary 2, equation (1.9)]
{montgomeryvaughan} states
\[
 \left|\sum_{j\ne k}\frac{a_j\overline{a_k}}{\omega_j-\omega_k}\right|
 \le\frac{3\pi}{2}\sum_j\frac{|a_j|^2}{\delta_j}.
\]
Integrating the exponential sum and applying this inequality to its two
endpoint phase products gives
\begin{equation}\label{lj:hilbertaverage}
 \left|\frac1H\int_0^H\left|\sum_ja_je^{i\omega_jy}\right|^2\dd y
                     -\sum_j|a_j|^2\right|
 \le\frac{3\pi}{H}\sum_j\frac{|a_j|^2}{\delta_j}.
\end{equation}
The local weighted gaps, rather than a common minimum gap, are essential.

Temporarily omit only \(n=N\), freeze the center data in
\eqref{lj:exactdata}, and write
\[
 g_-(y)=\sum_{n<N}w_n
       \{r_n\sin(\phi_n-\lambda_ny)-c_n\cos(\phi_n-\lambda_ny)\}.
\]
Its reflected exponential coefficients are exactly
\[
 a_{n,-}=\frac{w_n}{2}e^{i\phi_n}(-c_n-ir_n),\
 \omega_{n,-}=-\lambda_n,\qquad
 a_{n,+}=\overline{a_{n,-}},\
 \omega_{n,+}=\lambda_n.
\]
Its diagonal is therefore
\begin{equation}\label{lj:reflecteddiagonal}
 \sum_{n<N}(|a_{n,-}|^2+|a_{n,+}|^2)
             =\frac12\sum_{n<N}w_n^2(r_n^2+c_n^2).
\end{equation}
This frequency set retains the difference phases and the reflected sum
phases, including the oscillating reflected diagonal.

The exact phase center has the expansion
\begin{equation}\label{lj:phasecenter}
 M_\phi=e^{\Omega/c}
       =\sqrt{x/(4\pi)+t/16}+O(x^{-3/2}),\qquad c=\tfrac12+o(1).
\end{equation}
One may verify the needed accuracy directly from
\[
 \alpha_r=L/2+\tfrac14\log(1+x^{-2})-(1+x^2)^{-1},\qquad
 \alpha_i=3x/(1+x^2)-\tfrac12\arctan x,
\]
\[
 U=(7x^2-5)/(1+x^2)^2,\qquad V=x(x^2+5)/(1+x^2)^2.
\]
In fact \(\Omega/c=\alpha_r-t\alpha_iV/(2(1+tU/2))
=L/2+t\pi/(8x)+O(x^{-2})\);
the logarithm of the square root in \eqref{lj:phasecenter} has the same
first two terms. Hence, for \(n<N\) and small \(t\),
\[
 \lambda_n=c\log(M_\phi/n)\ge.45\log(N/n)>0.
\]
The possible error \(O(x^{-2})\) in \(\log(M_\phi/N)\) is negligible
beside \(\log(N/(N-1))\ge1/N\).
Same-sign successor and predecessor gaps are at least \(c/(2n)\)
and \(c/n\). An opposite-sign gap is at least \(\lambda_n\);
concavity of \(u\log(N/u)\) and its endpoint values on \([1,N-1]\)
give \(n\log(N/n)\ge1/2\). Thus every local gap obeys
\begin{equation}\label{lj:localgaps}
 \delta_{n,\pm}\ge1/(8n)\quad(n<N).
\end{equation}
No division by the possibly very small or slightly negative
\(\lambda_N\) has occurred.

The complete positive budget has the sharper uniform asymptotic
\begin{equation}\label{lj:derivativebudget}
 B_d:=\sum_{n\le N}nw_n^2(r_n^2+c_n^2)
       =\left(\frac{8t^2}{\kappa^2}+o(t^2)\right)e^{2\mathfrak a/t}.
\end{equation}
To see this, put \(u=e^{L/2-v}\) in the index integral. The exact
weight gives
\[
 uw(u)^2\dd u
 =e^{2\mathfrak a/t}e^{-v+tv^2/2}(1+o(1))\dd v,
 \qquad r(u)=2tv/\kappa+o(t),\quad c(u)=O(t/x)
\]
on bounded \(v\)-intervals. Since \(0\le tv\le1\), the normalized
integrand is dominated by
\(C t^2(1+v)^2e^{-v/2}\). The errors in \(r/t\) tend uniformly to
zero, and the \(c(u)^2\) contribution is bounded by
\(O(t^2x^{-2}\sum nw_n^2)\).
On each bounded \(v\)-interval the mesh tends exponentially to zero.
The discrete-integral error is at most a polynomial in \(L\):
indeed \(uw(u)^2=O(1)\) on the full index range, and its variation
and that of \(r(u)^2+c(u)^2\) are polynomially bounded in \(L\).
This error is negligible beside \(t^2e^{2\mathfrak a/t}\).
The floor and the interval beyond \(e^{L/2}\) are negligible as well.
Dominated convergence now gives
\[
 \frac{4t^2}{\kappa^2}\int_0^\infty v^2e^{-v}\dd v
                         =\frac{8t^2}{\kappa^2},
\]
proving \eqref{lj:derivativebudget}.

Equations \eqref{lj:hilbertaverage}--\eqref{lj:localgaps} yield
\begin{equation}\label{lj:frozenpaid}
 \frac1H\int_0^Hg_-(y)^2\dd y
 \ge\frac12\sum_{n<N}w_n^2(r_n^2+c_n^2)
             -\frac{12\pi}{H}\sum_{n<N}nw_n^2(r_n^2+c_n^2).
\end{equation}
The factor \(12\pi\) includes both reflected coefficient families.
Since \(w_1=1,c_1=0,r_1=1+o(1)\), the diagonal is at least
\(1/2+o(1)\). Taking \(H=H_d\) and then choosing an absolute
\(D\) sufficiently large makes the right side at least \(3/8\).
Restore the exact last frozen term, obtaining \(g_0\). Its supremum is
\[
 w_N\sqrt{r_N^2+c_N^2}=o(1),
\]
because \(w_N=O(e^{-\mathfrak b/t})\) and \(r_N,c_N\) are bounded.
The normalized \(L^2\) triangle inequality restores this coefficient
with an \(o(1)\) loss.

It remains to pay genuine physical movement. The averaging range satisfies
\begin{equation}\label{lj:physicalsize}
 H_d/N=O(t^2e^{-\kappa^2/(8t)})=o(1),\qquad
 H_d^2/x=O(t^4e^{-\kappa^2/(4t)})=o(1).
\end{equation}
At fixed time, uniformly for \(n\le N\) and \(0\le y\le H_d\),
\[
 |(\log w_n)'|=O(x^{-1}),\quad
 |\phi_n''|=O(x^{-1}),\quad
 |r_n'|=O((Lx)^{-1}),\quad |c_n'|=O(t/x^2).
\]
Here the last two coefficients use the fixed denominator \(L\).
Using \(|r_n|+|c_n|\le2\) and
\(\sum_{n\le N}w_n=O(e^{\mathfrak a/t})\), proved below in
\eqref{lj:absolutejets}, Taylor's formula in each exact summand gives
\begin{equation}\label{lj:movement}
 \sup_{0\le y\le H_d}|G_F(y)-g_0(y)|
 \le C e^{\mathfrak a/t}(H_d/x+H_d^2/x)=o(1).
\end{equation}
This pays amplitude, derivative-weight, and nonlinear phase changes.
The largest exponential factor is
\(e^{(5\mathfrak a-\kappa)/t}\), with
\(5\mathfrak a-\kappa=\kappa(4-5\kappa)/16\le-1/16\).
The first inequality in \eqref{lj:averagelower} follows.

For the exact heat function use, at each center \(x+y\), a separate
disk of radius \(1/L_y\), where \(L_y=\log((x+y)/(4\pi))\).
By \eqref{lj:physicalsize}, its continuous natural cutoff and that of
the whole interval vary by \(o(1)\); every natural integer cutoff
differs from the fixed \(N\) by at most one. The preceding complete
disk majorant pays this coefficient, normalization, and lower-half
reflection before Cauchy's estimate is applied. Thus
\[
 |Q_t'(x+y)-F_{t,N}'(x+y)|\le L_y\eta_{N,y}.
\]
The strict reserves of that ledger persist for
\(tL_y\in[1,2+o(1)]\), in particular \(tL_y\le2.01\) for small \(t\).
Also \(L_y/L=1+o(1)\), and the change in its error exponent is
\(O(H_d/x)=o(1)\). Consequently
\begin{equation}\label{lj:physicalcauchy}
 \sup_{0\le y\le H_d}|G_Q(y)-G_F(y)|
 \le\sup_{0\le y\le H_d}2(L_y/L)\eta_{N,y}
                  =O(e^{-\mathfrak b/t})=o(1).
\end{equation}
This uses separate small disks, not an unpaid remainder on a disk of
radius \(H_d\). A final normalized \(L^2\) triangle inequality proves
the second assertion.
\end{proof}

For comparison, the same reflected Hilbert calculation for the complete
value and derivative coordinates has budget
\(\sum nw_n^2(1+r_n^2+c_n^2)=O(e^{2\mathfrak a/t})\).
It removes the earlier logarithm already on the sufficient range
\(H=De^{2\mathfrak a/t}\). The derivative theorem is shorter because
its endpoint squared weight contains \(r_n^2\asymp t^2v^2\).
Neither sufficient range is a necessary lower bound for all possible
probe methods.

\begin{remark}[The remaining Taylor deficit]\label{lj:taylor}
The complete absolute jet and Cauchy budgets give
\(|Q_t''|\le M_2=C_2e^{\mathfrak a/t}\) on the physical interval.
If its left endpoint is a collision, Taylor's theorem gives
\[
 |Q_t'(x+y)|\le M_2y,\qquad
 \frac1H\int_0^H(2Q_t'(x+y)/L)^2\dd y
                         \le\frac{4M_2^2H^2}{3L^2}.
\]
This proved upper bound is small only on the sufficient Taylor scale
\(H=O(Le^{-\mathfrak a/t})\), whereas
\[
 \frac{H_d}{Le^{-\mathfrak a/t}}
               =\frac{Dt^2}{L}e^{3\mathfrak a/t}\longrightarrow\infty.
\]
At \(H=H_d\) the lower average only forces
\(\sup|Q_t''|\ge\sqrt{3/32}\,L/H_d\), compatible with the
absolute upper budget. The improved average therefore does not exclude
an isolated collision.
It also yields a limited measure statement: the absolute first-jet
budget gives \(\sup|G_Q|\le Ct e^{\mathfrak a/t}\), so, for
\(E=\{y\in[0,H_d]:|G_Q(y)|\ge1/4\}\),
\[
 |E|/H_d\ge c_0t^{-2}e^{-2\mathfrak a/t},\qquad |E|\ge c_1>0.
\]
Indeed its mean square is at most
\(1/16+C^2t^2e^{2\mathfrak a/t}|E|/H_d\).
The lower bound on relative measure tends to zero and still permits
isolated double zeros.
\end{remark}

\section{Threshold collision jets and a paid Laguerre condition}
\label{lj:threshold}

The heat equation supplies exact identities at a collision. The extra
sign used below instead comes from the global all-real property at the
threshold time. These two inputs must be kept distinct.

\subsection{Normalized backward evolution and collision compatibility}

Let \(\ell(t,x)=\log A_t(x)\),
\(b_t=\partial_x\ell\), \(a_0=\partial_t\ell\), and
\(\mathcal V_t=a_0+\partial_xb_t+b_t^2\). Direct differentiation of
\(H_t=A_tQ_t\) in \eqref{eq:heat} gives
\begin{equation}\label{lj:normalizedpde}
 \partial_tQ_t=-Q_t''-2b_tQ_t'-\mathcal V_tQ_t,\qquad
 (\partial_t+a_0)Q_t=-(\partial_x+b_t)^2Q_t.
\end{equation}
At \(Q_t=Q_t'=0\),
\begin{equation}\label{lj:collisionpde}
 \partial_tQ_t=-Q_t'',\qquad
 \partial_tQ_t'=-Q_t'''-2b_tQ_t''.
\end{equation}
At an ordinary double zero, the Jacobian of \((Q_t,Q_t')\) in
column order \((x,t)\) therefore has determinant \((Q_t'')^2>0\).
The normalizer and its transport coefficient retain the transverse
ordinary collision geometry.

The exact normalizer coefficients, with \(s=(1-ix)/2\), are
\begin{align}
 a_0&=\frac14\Rea(\alpha(s)^2),&
 b_t&=\frac12\Ima m_t'(s),&
 m_t'&=\alpha(1+t\alpha'/2),\nonumber\\
 \partial_xb_t
   &=-\frac14\Rea\left\{\alpha'
        +\frac t2(\alpha'^2+\alpha\alpha'')\right\},&
 \alpha''(s)&=s^{-3}+2(s-1)^{-3}-\frac12s^{-2}.
 \label{lj:bx}
\end{align}
In \eqref{sc:scaling}, \(\partial_xb_t=O(x^{-2})\):
\(\Rea\alpha'=O(x^{-2})\),
\(\alpha'^2=O(x^{-2})\), and the remaining product after
multiplication by \(t\) is \(O(x^{-2})\), since \(t\log x=O(1)\).
The heat correction in the normalizer is included.

There is an exact carrier conjugation as well. Put
\(M=M_t(s)=A_te^{i\theta_t}\) and
\(\mathcal U=H_t/M=e^{-i\theta_t}Q_t\). Since
\(M_x/M=b_t-i\Omega\),
\begin{equation}\label{lj:carrierpde}
 \partial_t\mathcal U=-\mathcal U''-2(b_t-i\Omega)\mathcal U'
                       -v_M\mathcal U,\qquad
 v_M=\frac14\{\alpha^2-(m_t')^2-m_t''\},
\end{equation}
where \(m_t''=\alpha'+t(\alpha'^2+\alpha\alpha'')/2\).
The real potential in \eqref{lj:normalizedpde} satisfies
\(\mathcal V_t=\Omega^2+\Rea v_M\).
At \(\mathcal U=\mathcal U'=0\), this still gives
\(\partial_t\mathcal U=-\mathcal U''\); conjugation alone supplies
no nonvanishing condition.

For the genuine unnormalized function define
\(W=(H_t')^2-H_tH_t''\). Its backward evolution is
\begin{equation}\label{lj:laguerreevolution}
 \partial_tW=-W''+2\{(H_t'')^2-H_t'H_t'''\}.
\end{equation}
At an ordinary double real collision \((T,x_*)\),
\[
 W=0,\qquad W''=\partial_tW=(H_T''(x_*))^2>0.
\]
In particular \(W(T-h,x_*)=-h(H_T''(x_*))^2+O(h^2)<0\)
for small \(h>0\), consistent with nonreal roots immediately before
the collision. Equation \eqref{lj:laguerreevolution} is backward
parabolic with a nonzero collision source. All-real Laguerre
nonnegativity at time \(T\) does not turn it into an exclusion by a
forward maximum principle.
The corresponding exact normalized identity is
\begin{equation}\label{lj:normalizedlaguerre}
 W/A_t^2=(Q_t')^2-Q_tQ_t''-(\partial_xb_t)Q_t^2.
\end{equation}
At an ordinary double collision its value and first spatial derivative
vanish, while its time and second spatial derivatives equal
\((Q_t'')^2\). These identities are proved by differentiating
\(H_t=A_tQ_t\) and \(\partial_tH_t=-H_t''\); no signed arithmetic
estimate has entered.

\subsection{The deflated mirror inequality and its multiplicity hierarchy}

If \(\Lambda>0\), Proposition~\ref{prop:finitecollision} supplies a
finite multiple real zero at time \(\Lambda\). At that specific time
all zeros are real by the threshold theorem. This global fact is not
inferred from a maximizing root, the local PDE, or an arbitrary real
collision at positive time.

We use the canonical-product facts in
\cite[Section 3 and the proof of Proposition 3.1]{polymath}:
\(H_T\) is real entire of order one, even, nonzero at zero, and has
a locally uniformly convergent Hadamard product grouped in opposite-root
pairs. At an all-real time it is a constant times
\(\prod_{\rho>0}(1-z^2/\rho^2)\), with multiplicities.
The inverse-square sum obtained by a second logarithmic derivative
converges absolutely off the roots. Deflating finitely many roots
removes exactly their terms from that sum; a possible linear exponential
factor has zero second logarithmic derivative.

\begin{proposition}[Deflated forced-mirror condition]\label{lj:mirror}
Suppose all zeros of \(H_T\) are real, and \(x_*>0\) is a root of
multiplicity \(m\ge2\). Put \(H_j=H_T^{(j)}(x_*)\). Then
\begin{equation}\label{lj:unnormalizedmirror}
 2H_3^2-3H_2H_4\ge9H_2^2/x_*^2.
\end{equation}
Writing \(q_j=Q_T^{(j)}(x_*)\) and
\begin{equation}\label{lj:gamma}
 \gamma(T,x_*)=18\partial_xb_T(x_*)+9/x_*^2,\qquad
 \mathscr L(q)=2q_3^2-3q_2q_4-\gamma q_2^2,
\end{equation}
the exact normalized condition is
\begin{equation}\label{lj:normalizedmirror}
 \mathscr L(q)\ge0.
\end{equation}
It applies in particular to every finite positive threshold collision.
For \(m\ge4\) this fourth-jet condition is the uninformative equality
zero.
\end{proposition}

\begin{proof}
Deflate two copies, \(g(z)=H_T(z)/(z-x_*)^2\). Evenness leaves at
least two copies of the real mirror root \(-x_*\).
For real \(z\) away from roots, the canonical product gives
\[
 (g')^2-gg''
       =g(z)^2\sum_{\rho\ {\rm root\ of}\ g}\frac1{(z-\rho)^2}
       \ge\frac{2g(z)^2}{(z+x_*)^2}.
\]
Pass to the continuous limit \(z\to x_*\), along nonroots. This is
valid also when \(m>2\), for which \(g(x_*)=0\); a logarithmic
derivative is not taken at that zero. The Taylor coefficients are
\[
 g(x_*)=H_2/2,\qquad g'(x_*)=H_3/6,\qquad g''(x_*)=H_4/12.
\]
Multiplication by 72 proves \eqref{lj:unnormalizedmirror}.
At the collision \(q_0=q_1=0\), and exact product differentiation gives
\[
 H_2=A_Tq_2,\quad
 H_3=A_T(q_3+3b_Tq_2),\quad
 H_4=A_T\{q_4+4b_Tq_3+6(\partial_xb_T+b_T^2)q_2\}.
\]
Substitution cancels all \(b_T\)-terms except
\(-18(\partial_xb_T)q_2^2\), proving
\eqref{lj:normalizedmirror}.
At exact multiplicity three \(q_2=0,q_3\ne0\); at multiplicity
at least four \(q_2=q_3=0\). The latter case is allowed, not excluded,
by this condition.
\end{proof}

At an exact multiplicity \(m\ge2\), deflating all \(m\) copies instead
gives the higher-order condition
\begin{equation}\label{lj:hierarchy}
 \frac{q_{m+1}^2}{(m+1)^2}
 -\frac{2q_mq_{m+2}}{(m+1)(m+2)}
 -\left(\partial_xb_T+\frac{m}{4x_*^2}\right)q_m^2\ge0.
\end{equation}
To derive it, use \(g=H_T/(z-x_*)^m\), which is nonzero at \(x_*\)
and retains \(m\) mirror copies. Its inverse-square sum is at least
\(m/(4x_*^2)\). Substituting
\[
 g(x_*)=H_m/m!,\quad
 g'(x_*)=H_{m+1}/(m+1)!,\quad
 g''(x_*)=2H_{m+2}/(m+2)!
\]
and multiplying by \((m!)^2\) gives the unnormalized version.
At exact multiplicity \(m\),
\begin{align*}
 H_m&=A_Tq_m,\\
 H_{m+1}&=A_T\{q_{m+1}+(m+1)b_Tq_m\},\\
 H_{m+2}&=A_T\{q_{m+2}+(m+2)b_Tq_{m+1}
               +\tbinom{m+2}{2}(\partial_xb_T+b_T^2)q_m\}.
\end{align*}
The \(b_T\)-terms cancel, leaving \eqref{lj:hierarchy}.
This hierarchy requires jets \(m,m+1,m+2\) and their own approximation
payments. It is not resolved by a bounded second-through-fourth-jet
certificate.

\subsection{Complete measured payment for the threshold jet}

The sign condition \eqref{lj:normalizedmirror} becomes an arithmetic
target only after all products of approximation errors are paid.
For an integer \(K\le N\), fixed when differentiating, put
\(f_j=F_{t,K}^{(j)}(x)\). Choose proved nonnegative tail payments
\(\tau_j\ge|F_{t,N}^{(j)}-F_{t,K}^{(j)}|\), \(j=2,3,4\), and define
\begin{equation}\label{lj:jeterrors}
 \delta_j=j!L^j\eta_N+\tau_j.
\end{equation}
Then \(|Q_t^{(j)}-f_j|\le\delta_j\) by the original disk payment
followed by exact signed subtraction.

\begin{proposition}[Measured paid threshold test]\label{lj:paidtest}
At each parameter point in \eqref{sc:scaling}, define
\begin{align}
 \Delta(f,\delta)={}&4|f_3|\delta_3+2\delta_3^2
       +3|f_2|\delta_4+3|f_4|\delta_2+3\delta_2\delta_4\nonumber\\
       &+|\gamma|\{2|f_2|\delta_2+\delta_2^2\},
       \qquad \gamma=18\partial_xb_t+9/x^2.
       \label{lj:measuredpayment}
\end{align}
The strict genuine arithmetic inequality
\begin{equation}\label{lj:negativetest}
 2f_3^2-3f_2f_4-\gamma f_2^2<-\Delta(f,\delta)
\end{equation}
excludes an all-real-time collision at that point.
The full cutoff \(K=N\) has \(\Delta_N=o(1)\) uniformly.
The original value core \(K_v\) also has a uniformly vanishing absolute
payment, in fact
\(\Delta_{K_v}=O(N^{-1/6})\). The deeper derivative core \(K_d\)
does not inherit this conclusion uniformly from the displayed absolute
jet bounds.
\end{proposition}

\begin{proof}
Write \(q_j=f_j+e_j\), \(|e_j|\le\delta_j\).
Expansion of \(2q_3^2\), \(-3q_2q_4\), and \(-\gamma q_2^2\)
gives
\begin{equation}\label{lj:productpayment}
 |\mathscr L(q)-\mathscr L(f)|\le\Delta(f,\delta).
\end{equation}
In particular the \(3\delta_2\delta_4\) and
\(|\gamma|\delta_2^2\) products cannot be discarded.
At an all-real-time collision Proposition~\ref{lj:mirror} gives
\(\mathscr L(q)\ge0\), contradicting
\eqref{lj:negativetest} and \eqref{lj:productpayment}.

For the absolute ledger, for every fixed \(0\le j\le4\) and either
retained prefix,
\begin{equation}\label{lj:absolutejets}
 |F_{t,N}^{(j)}|+|F_{t,K}^{(j)}|
       =O_j(e^{\mathfrak a/t})
       =O_j(N^{(4-\kappa)/8}).
\end{equation}
To prove this without signed cancellation, put \(y=L/2-\log u\).
The index mass satisfies
\[
 uw(u)\le C e^{\mathfrak a/t}e^{-y/4}
                       \quad(0\le y\le L/2),
\]
since its exponent is
\(\mathfrak a/t-y/2+ty^2/4+O(x^{-2})\) and
\(ty^2/4\le\kappa y/8\le y/4\).
The complete raw derivative multiplier has size
\(O_j((1+y)^j)\), by \eqref{lj:rawmultiplier} and its fixed-time
spatial derivatives. Index integration is therefore
\(O_j(e^{\mathfrak a/t})\).
For a direct discrete comparison, \(w(u)\) decreases throughout this
range, as does \((1+L/2-\log u)^j\); their product is decreasing.
The initial term is \(O_j(L^j)\), absorbed by \(e^{\mathfrak a/t}\).
The negligible endpoint displacement beyond \(e^{L/2}\) is absorbed
as well. This proves \eqref{lj:absolutejets}.
The same proof on the physical interval, followed by the separate
small-disk Cauchy estimates, supplies the curvature and first-jet
budgets used in Remark~\ref{lj:taylor}.

For \(K=N\), \(\tau_j=0\). Equations
\eqref{lj:cauchypayment}, \eqref{lj:bx}, and
\eqref{lj:measuredpayment} yield
\begin{equation}\label{lj:fullquadraticpayment}
 \Delta_N=O\left(L^4e^{(\mathfrak a-\mathfrak b)/t}
                         +L^6e^{-2\mathfrak b/t}\right)
         =O\left(L^4e^{-\kappa^2/(8t)}
                         +L^6e^{-2\mathfrak b/t}\right)=o(1).
\end{equation}
Thus growing absolute jets do not prevent a vanishing absolute error
for this quadratic test.

For \(K=K_v\), use \eqref{lj:shortjets} in \eqref{lj:jeterrors}.
The largest linear tail product is the second-derivative payment times
an absolute fourth jet:
\[
 e^{\mathfrak a/t}\tau_2
 =O\left(N^{7/12-s_\kappa-1/2+(4-\kappa)/8}\right)
 =O(N^{1/12-\kappa/4})=O(N^{-1/6}).
\]
The other linear products have extra powers \(N^{-1/4}\) and
\(N^{-1/2}\); quadratic tail products decay faster, and
\(\gamma=O(x^{-2})\). Consequently
\begin{equation}\label{lj:corequadraticpayment}
 \Delta_{K_v}
 =O\left(N^{-1/6}+L^4e^{-\kappa^2/(8t)}
                         +L^6e^{-2\mathfrak b/t}\right)
 =O(N^{-1/6}).
\end{equation}
The last equality uses the strict reserve
\(e^{-\kappa^2/(8t)}=O(N^{-\kappa/4})\), which absorbs every
fixed power of \(L\) relative to \(N^{-1/6}\).

For \(K=K_d\), \eqref{lj:deepjets} pays every individual derivative
tail, but the same absolute linear products have powers
\[
 N^{11/24-\kappa/4},\qquad
 N^{1/3-\kappa/4},\qquad N^{5/24-\kappa/4}.
\]
The first budget is positive for \(\kappa<11/6\).
It tends to zero only on closed subranges
\([11/6+\varepsilon,2]\); at \(\kappa=11/6\) it is merely bounded.
This is a limitation of the absolute transfer, not evidence that the
genuine quadratic difference grows. The measured test
\eqref{lj:measuredpayment}--\eqref{lj:negativetest} remains valid
with these payments retained.
\end{proof}

No negative sign in \eqref{lj:negativetest} is proved for the actual
height phases. At a putative positive threshold collision, one may
condition that arithmetic task on the necessary small coordinates
\eqref{lj:mixedcollision}; those conditions alone have not been shown
to force the opposite sign. The all-real hypothesis is still required.
The test is not an assertion about an arbitrary positive-time collision
when other roots are nonreal, and this shrinking-time regime leaves
other times and heights untreated.

\subsection{A sharp model for the PDE and Laguerre mechanism}

The mirror coefficient and compatibility with a positive threshold are
sharp in a general even polynomial heat-flow class.

\begin{proposition}[Quartic threshold model]\label{lj:quartic}
Fix \(T>0\), \(a>0\), set \(\delta=t-T\), and put
\begin{equation}\label{lj:quarticflow}
 P_t(z)=z^4-(2a^2+12\delta)z^2
          +a^4+4a^2\delta+12\delta^2
       =e^{-\delta\partial_z^2}(z^2-a^2)^2.
\end{equation}
Then \(\partial_tP_t=-P_t''\), and all its zeros are real exactly
for \(t\ge T\). At \((T,a)\) it attains equality in
\eqref{lj:unnormalizedmirror}. If \(a^2>6T\), it also satisfies
\(P_t(iy)>0\) for every \(t\ge0\) and \(y\in\mathbb R\).
\end{proposition}

\begin{proof}
The finite exponential in \eqref{lj:quarticflow}, or direct
differentiation, proves the heat equation. As a quadratic in \(z^2\),
its discriminant is \(32\delta(a^2+3\delta)\), and its constant
coefficient is positive for every real \(\delta\).
For \(\delta\ge0\), both roots in \(z^2\) are positive.
For \(-a^2/3<\delta<0\), the discriminant is negative.
For \(\delta\le-a^2/3\), both roots in \(z^2\) are negative,
including the equal-root endpoint. Thus the all-real threshold is
exactly \(T\).
If \(a^2>6T\), substituting \(z=iy\) leaves positive coefficients
in \(y^4\), \(y^2\), and the constant for all \(t\ge0\).
At the collision,
\[
 P_2=8a^2,\qquad P_3=24a,\qquad P_4=24,\qquad
 2P_3^2-3P_2P_4=576a^2=9P_2^2/a^2.
\]
This is the asserted saturation.
\end{proof}

Dividing the model by the same explicit zero-free normalizer on its
analytic domain yields \eqref{lj:normalizedpde}, with the same
normalizer coefficients, and retains its positive-time collisions.
The model therefore shows that the backward heat PDE, normalization,
threshold real-rootedness, and the stated Laguerre signs alone do not
exclude an ordinary positive collision. It is not the zeta heat
function, has no zeta infinite zero set or arithmetic coefficients,
and is not asserted to come from a positive theta kernel.
Additional genuine signed arithmetic could distinguish it.
The remaining task is a proved local joint lower bound, or a strict
paid opposite jet test for the actual phases, with a higher-jet
hierarchy where necessary.


\section{The next analytic route: paid collision exclusion}
\label{sec:outlook}
The preceding sections supply signed reductions and complete error
payments in addition to the geometric and landing tools. They do not
exclude a genuine positive-time collision. Corollary~\ref{cor:rhcollision}
permits the direct target $H_t=H_t'=0$ at every positive time. For an
RH argument it would also suffice to exclude positive all-real threshold
collisions: if $\Lambda>0$, Proposition~\ref{prop:finitecollision}
produces exactly such an event. The latter route may use the extra
real-root sign in the threshold jet test. Neither route has been completed,
and the shrinking sector alone leaves other heights and scaling regimes.

\subsection{A conditional amplitude and phase criterion}
On the real axis write the fixed-cutoff approximation as
$F=2\Rea(e^{i\theta}P)$. Wherever $P\ne0$, put
$m=|P|$ and $\phi=\theta+\arg P$ locally. Then
\begin{equation}\label{eq:phaseform}
 F=2m\cos\phi,\qquad
 F'=2m'\cos\phi-2m\phi'\sin\phi,
 \qquad\phi'=\theta'+\Ima(P'/P).
\end{equation}
No absolute choice of phase branch enters the derivative.

\begin{lemma}[Error-paid phase exclusion]\label{lem:phase}
Suppose $|Q-F|\le\epsilon_0$ and $|Q'-F'|\le\epsilon_1$ on a
region, and
$m\ge m_0>\epsilon_0/2$, $|m'|\le m_1$, and $|\phi'|\ge\Omega>0$.
Then
\begin{equation}\label{eq:phasecriterion}
 2\Omega\sqrt{m_0^2-\epsilon_0^2/4}
       -\frac{m_1}{m_0}\epsilon_0>\epsilon_1
\end{equation}
excludes a common zero of $Q,Q'$.
A simpler sufficient criterion, under
$|P|\ge p_0$, $|P'|\le p_1$, $|\theta'|\ge\omega_0>0$, is
\begin{equation}\label{eq:carriercriterion}
 \epsilon_0<2p_0,\qquad
 2\omega_0\sqrt{p_0^2-\epsilon_0^2/4}>2p_1+\epsilon_1.
\end{equation}
\end{lemma}
\begin{proof}
At a candidate common zero, $|F|\le\epsilon_0$ and
$|F'|\le\epsilon_1$. The first bound gives
$m|\sin\phi|\ge\sqrt{m^2-\epsilon_0^2/4}$, while the amplitude
derivative contribution in \eqref{eq:phaseform} is at most
$(m_1/m_0)\epsilon_0$. The reverse triangle inequality now
contradicts \eqref{eq:phasecriterion}. For the second version,
differentiate $F=2\Rea(e^{i\theta}P)$ directly. The carrier term
has modulus at least
$2\omega_0\sqrt{p_0^2-\epsilon_0^2/4}$, and the remaining
term is at most $2p_1$.
\end{proof}
These are conditional comparison lemmas. No many-term amplitude
lower bound or phase-speed bound has been proved here. A global
nonvanishing requirement on $P$ may be stronger than necessary:
a successful argument may instead use the two candidate collision
conditions jointly. The errors must come from a proved complex
neighborhood estimate, as in Section~\ref{sec:normalization}.

\subsection{The shrinking-time regime and an absolute-value obstruction}
The proved signed reductions concern the genuine many-term sum in
\begin{equation}\label{eq:kapparegime}
 \kappa=t\log(x/(4\pi))\in[1,2],\qquad t\downarrow0.
\end{equation}
The full error, growing-edge cancellation, and shorter derivative
average have now been proved in this sector. A uniform local signed
implication at candidate collisions is still missing. Proposition~\ref{prop:absmass} explains why the successful short-edge
subtraction cannot simply be extended by positive tail summation.
Decreasing the constant in Corollary~\ref{cor:eventual} does not
settle the shrinking-time target.

\subsection{The two remaining local signed targets}
At a candidate collision the value and derivative cores satisfy the
small-coordinate bounds proved above. One sufficient next input is a
uniform lower bound for that genuine mixed-cutoff vector exceeding its
complete error budget. Because its two cutoffs differ, its second
coordinate is not the derivative of its first. A same-cutoff formulation
may instead retain the original value core and its paid derivative.

At an all-real candidate there is a second interface. Evaluate the
fourth-derivative expression $\mathscr L$ on the full approximant or the
value core, and retain the measured product error $\Delta$. A proof that
the candidate value and derivative conditions force
$\mathscr L(F)<-\Delta$ would contradict the necessary threshold sign.
The approximation payment is already available; the strict signed
implication has not been proved. A collection of small individual
derivative errors is insufficient unless their products are also paid.
The deeper derivative core has no uniform vanishing product budget on
the entire sector. Multiplicity at least four makes the fourth-derivative
condition vacuous, so a successful local argument must also handle that
case, either directly or through the higher deflation hierarchy.

Complete multiplicative twists with freely chosen prime phases satisfy
less structure than the physical height orbit. Their joint zeros prevent
a lower bound based only on complete multiplicativity. The shorter
derivative average likewise permits isolated double zeros: its interval
still exceeds the scale on which the paid curvature bound transfers a
collision to a small derivative. Thus the next estimate must use genuine
arithmetic correlations at the candidate, or a stronger local mechanism.

\subsection{Other analytic tools and their proof checkpoints}
The translated-probe discussion in \cite[Section~8]{planat} motivates
retaining correlated information from the same arithmetic sum.
The complete derivative average in Theorem~\ref{lj:derivativeaverage}
now pays such an averaged estimate, while Proposition~\ref{sc:packet}
shows why a fixed number of bounded packet probes cannot supply diagonal
coercivity. A next useful local estimate must use the candidate conditions
and retain the adverse signed cross terms and physical errors.
Generic derivative nondegeneracy alone permits the transverse collision
\eqref{eq:jacobian}.

The parallel arithmetic investigations in the repository suggest
three other proof checkpoints. First, a fixed-scale contraction of
the actual normalized prime variance would be an endpoint
mechanism if it bounded the entire net energy change, including
the squared defect and every boundary term, by a strict negative
fraction plus a power remainder. A negative covariance alone does
not establish that budget. Second, short-family extraction can
approach the critical line only under a hierarchy of actual
moment and compensation estimates whose positive scale parameters
approach zero; the retained deletion variance must be paid.
Third, the developed gated mixed covariance is an arithmetic
target for a first local strip improvement, but no iteration of
that architecture to RH has been proved. The detailed conditional
hypotheses and stopping conditions remain in the companion
short-family and mixed-character manuscripts and the continuation
note dated 9 October 2026. These programs are alternatives to be
judged by signed arithmetic estimates, not by equivalent endpoint
reformulations or expanding numerical campaigns.

For the present heat route, the next bounded checkpoint is the
candidate-conditioned signed jet implication or a direct local joint
lower bound in \eqref{eq:kapparegime}, including higher multiplicities
and every approximation payment. Its proof in this sector would be a
partial exclusion result; a separate coverage argument would still be
needed for RH. The collective attraction tools remain available for a
positive-time application, while an additional finite field improvement
alone remains compatible with a positive collision.

\appendix
% Appendix fragment. Prepared for Edward Baker, 9 October 2026.
% Model: gpt-6.1-sol (Codex); reasoning effort: ultra.
% Same-model assistance and cross-review are not independent validation.

\section{An effective probe and its spatial cost}\label{app:effective-probe}

This appendix records an effective consequence of the approximation theorem
of \cite[Theorem~1.3, equations (20)--(24)]{polymath}.  It uses analytic
tail bounds rather than evaluations of the heat function or a zero search.
The threshold is deliberately conservative.  The resulting attraction bound
applies throughout an unbounded spatial sector; the global consequence below
also needs a separate estimate for the complementary sector.

\begin{proposition}[Effective probe]\label{prop:effective-probe}
For \(x\ge 10^{16}\) and \(1/5\le t\le 3/10\),
\[
 H_t(x+9i/10)\ne0,\qquad
 L_{9/10}(t,x):=-\Im\frac{H_t'}{H_t}(x+9i/10)
 >\frac{12749}{2000}>6.
\]
Consequently, at any simple nonreal zero \(z=x+iy\) of globally maximal
imaginary height, with \(|x|\ge10^{16}\) and \(0<y\le1/5\), the exact
attraction field satisfies
\[
 E_t(z)\ge\frac{940}{231}>4.
\]
\end{proposition}

\begin{proof}
Put \(r=1/20\), \(h=x-r\), \(q_-=h/(4\pi)\),
\(\ell=\log q_-\), and
\[
 D_x=\{z:|z-(x+9i/10)|\le r\},\qquad
 N=\left\lfloor\sqrt{q_-+1/80}\right\rfloor.
\]
Every point \(z=u+iy\) of this disk satisfies
\(u>9\cdot10^{15}\) and \(17/20\le y\le19/20\), so the entire disk
lies in the upper-half-plane domain of the imported theorem.  There is no
need to reflect its error estimate below the real axis in this application.
The reflected Dirichlet term must still be included.  The inequalities
\(4\pi<13\), \(e^{67/2}<4\cdot10^{14}\), and
\(q_->9\cdot10^{15}/13>4\cdot10^{14}\) give
\(\ell>67/2\).  On the disk define
\[
 s=(1-iz)/2,\quad B_t(z)=M_t(s),\quad U_t(z)=H_t(z)/B_t(z),
\]
and retain the fixed holomorphic approximant
\[
 f_{t,N}(z)=P_{t,N}(s)+\gamma_t(z)P_{t,N}(1-s)=:P+R,
 \qquad \gamma_t=M_t(1-s)/M_t(s).
\]
The logarithm branches defining \(M_t\) are analytic here, and \(B_t\)
is nowhere zero.  In all the estimates that follow, \(N\) is fixed on
the whole disk and throughout the time interval.

The source coefficient bound has vanishing positive-part correction on
this disk, since \(1-3y+4y(1+y)/u^2<0\).  Thus
\[
 |p_{t,n}(s)|\le
 n^{-37/40-(t/4)(\ell-\log n)}.
\]
All retained indices obey \(\log n\le\ell/2+1/100\).
For \(2\le n\le16\), the weaker bound \(\log n<3\) gives
\(|p_{t,n}(s)|\le n^{-12/5}\); for \(n>16\), it gives
\(|p_{t,n}(s)|\le n^{-7/4}\).  Decreasing integral tails yield
\begin{align*}
 |P-1|&<\frac{19}{100}+\frac{19}{50}\frac57+\frac16
       =\frac{1319}{2100}<\frac23,\\
 \sum_{n=2}^{N}(\log n)|p_{t,n}(s)|
 &<\frac{13}{25}+\frac{13}{18}=\frac{559}{450}<\frac54.
\end{align*}
For completeness, the first head bound uses
\(2^{-12/5}<19/100\), \(2^{-7/5}<19/50\), and
\(\int_{16}^{\infty}u^{-7/4}\,du=1/6\).
The logarithm-weighted head is bounded by
\[
 \frac{19}{100}\frac7{10}
 +\frac{19}{50}\left(\frac12+\frac{25}{49}\right)
 =\frac{25327}{49000}<\frac{13}{25},
\]
using \(\log2<7/10\).  The logarithm-weighted tail is bounded by
\(\int_{16}^{\infty}(\log u)u^{-7/4}\,du<13/18\).
Directly from
\[
 \alpha'(s)=-\frac1{2s^2}-\frac1{(s-1)^2}+\frac1{2s},
\]
one has \(|\alpha'|\le2/h\) for both arguments.  Hence
\(|1+t\alpha'/2|<1001/1000\).  Differentiation of the fixed polynomial
gives
\[
 |P'|<\frac12\frac{1001}{1000}\frac54
       =\frac{1001}{1600}<\frac{63}{100}.
\]

For the reflected argument, \(\Re(1-s)\ge1/40\) and
\(\Re\alpha(1-s)\ge\ell/2-4/h^2\).  Its coefficients are bounded
by \(n^{-3/2}\) for \(n\le16\), and by \(n^{-17/20}\) for
\(n>16\).  The loss \(2t/h^2\) is covered by the strict slack in
these exponents.  Therefore
\[
 |P_{t,N}(1-s)|\le3+\frac{20}{3}N^{3/20}
                 \le3+7e^{3\ell/40}.
\]
The source reflection bound gives
\(|\gamma_t|\le(103/100)e^{-17\ell/40}\).  Consequently, uniformly
on the full disk,
\[
 |R|\le\frac{103}{100}
       \left(3e^{-17\ell/40}+7e^{-7\ell/20}\right)
       <\frac{7519}{10^8}<10^{-4}.
\]
Here \(e^{-1139/80}<10^{-6}\) and \(e^{-469/40}<10^{-5}\) suffice.
Cauchy's estimate for this holomorphic reflected term gives at the center
\[
 |R'|<20\frac{7519}{10^8}<\frac1{500}.
\]

We now pay the source remainder and every possible cutoff change before
applying Cauchy to the total remainder.  The natural cutoff at a source
point is
\(m=\lfloor\sqrt{u/(4\pi)+t/16}\rfloor\).
Across the disk and the time interval the squared parameters vary by
\(1/(40\pi)+1/160<1\).  Thus \(m\in\{N,N+1\}\).
Every index that actually occurs obeys
\[
 \log n\le\ell/2+\rho,\qquad
 \rho=\frac12\log\left(1+
            \frac{1/(40\pi)+3/160}{q_-}\right)<10^{-16}.
\]
In particular \(m>10^7+1\) and
\(\sum_{n\le m}|p_{t,n}(s)|\le7/3\).
The exact rewriting of the source's reflected term gives its modulus
multiplier as
\[
 |\gamma_t|m^{|\kappa|}n^y
 \le e^{.02y}
      \left(\frac{m^2}{u/(4\pi)}\right)^{y/2}
      e^{|\kappa|\log m}<\frac{103}{100}.
\]
This bound also holds for a crossing index \(n=N+1\).
Indeed the squared-index quotient is at most \(e^{2\rho}\), and
\(|\kappa|\log m<10^{-14}\), using
\(|\kappa|\le ty/[2(u-6)]\) and the monotonicity of \((\log u)/u\).
The total exponent is below \(.02\), and
\(e^{.02}\le50/49<103/100\).

For the source's two sum errors (equation (23) of \cite{polymath}),
the individual error exponent is bounded by
\[
 v(u)=\frac{.006(\log u)^2+1}{.99u}<1.1\cdot10^{-15}.
\]
This majorant decreases for \(u\ge9\cdot10^{15}\), and its endpoint
bound follows from \(\log(9\cdot10^{15})<37\).
Also \(|\log(u/(4\pi n^2))|\le\log u\); the possible negative
endpoint is bounded by \(\log(1+t/(16q))\), where \(q=u/(4\pi)\).
It follows from \(e^v-1\le v/(1-v)<1.12\cdot10^{-15}\), the
\(7/3\) sum bound, and the \(103/100\) reflected multiplier that
\(e_A+e_B<10^{-14}\).
For the source's remaining error (equation (24) of \cite{polymath}),
the two positive corrections have sum less than \(10^{-6}\): the first
is at most \(1.24(3+1)/(m-.125)<5\cdot10^{-7}\), and the second is at
most \((3\log u+16)/(.99u)<2\cdot10^{-14}\).
Its negative exponent has magnitude at least
\[
 \frac{37\ell}{80}+\frac{\ell^2}{80}
 \ge\frac{9447}{320}>29.5+10^{-6}.
\]
Thus \(e_{C,0}<e^{-29.5}<2\cdot10^{-13}\).

If \(m=N+1\), the complete difference between the two approximants is
\(p_{t,N+1}(s)+\gamma_t p_{t,N+1}(1-s)\).
For \(v=\log(N+1)\), one has \(\ell/2<v\le\ell/2+\rho\).
Completing the square \(v^2-\ell v=(v-\ell/2)^2-\ell^2/4\) gives
\[
 |p_{t,N+1}(s)|\le
 e^{-37\ell/80-\ell^2/80+\rho^2/20}
 <2\cdot10^{-13}.
\]
Its reflected companion is at most \(103/100\) times this value.
The entire cutoff jump is therefore less than \(5\cdot10^{-13}\).
All three payments yield
\[
 |U_t-f_{t,N}|<10^{-12}\quad\hbox{on }D_x,
 \qquad |(U_t-f_{t,N})'(x+9i/10)|<2\cdot10^{-11}.
\]
The latter is Cauchy's estimate with radius \(1/20\), applied to the
holomorphic fixed-cutoff remainder.  It does not differentiate a
moving-cutoff formula.

Finally, \(B_t'/B_t=-im_t'(s)/2\), where
\(m_t'\!=\alpha(1+t\alpha'/2)\).
The direct bounds
\[
 |\alpha|<3\ell/5,\qquad
 \Re\alpha\ge\ell/2-4/h^2,\qquad |\alpha'|\le2/h
\]
give \(\Re m_t'(s)\ge\ell/2-1/1000\).
For example \(\ell\le\log h\le\sqrt h\) and \(h>10^{15}\)
pay the heat correction.  Hence at the probe center
\[
 -\Im B_t'/B_t\ge67/8-1/2000,\qquad
 |U_t|>33/100,\qquad |U_t'|<13/20.
\]
In particular \(|U_t'/U_t|<2\), proving the stated nonvanishing and
\(L_{9/10}>67/8-1/2000-2=12749/2000\).
At a maximal zero with \(y\le1/5\), the ratio \((9/10)/y\) exceeds
\(\sqrt5\).  The exact-field comparison has factor one and gives
\[
 E_t(z)\ge\frac{L_{9/10}}{9/10}
               -\frac2{81/100-y^2}
 >\frac{20}{3}-\frac{200}{77}=\frac{940}{231}>4.
\]
Evenness gives the same assertion for negative real parts.
\end{proof}

The scalar reserves above are reproduced by
\href{../numerics/check_effective_collective_probe.py}{the effective-probe checker}
and retained in
\href{../numerics/effective_collective_probe_record_20261009.json}{its dated record}.
It verifies 105 exact rational, polynomial, and Taylor-series assertions,
including an outward enclosure of the conditional landing time below.
For a positive rational exponential argument \(a<122\), the degree-120
Taylor sum is a lower bound and adding
\( (a^{121}/121!)/(1-a/122)\) gives an upper bound.
The checker evaluates no heat functions, zeros, phases, or large cutoff
sums.  These scalar checks do not replace the imported approximation
theorem or the analytic arguments applying it.  Further proof details and
the scoped internal audit are in
\href{notes/7_EFFECTIVE_COLLECTIVE_PROBE_AND_SPATIAL_COST_20261009.md}{Heat Note 7}
and
\href{../reviews/EFFECTIVE_COLLECTIVE_PROBE_REVIEW_20261009.md}{its review}.

\begin{corollary}[Conditional sector landing]
Suppose the squared maximal height at \(t_0=1/5\) is at most \(1/25\),
and every globally maximizing nonreal zero remains in
\(|\Re z|\ge10^{16}\) until landing.  Then the flow has only real zeros by
\[
 t_{\mathrm{sector}}=\frac15+\frac1{16}\log\frac{33}{25}
 \in(0.21735,0.21736).
\]
\end{corollary}
\begin{proof}
On that population \(E>4\), so \(w'\le-2-16w\).
Integration, with the maximum-envelope continuation at switches and
multiple-root events, gives the stated time.
\end{proof}

The spatial premise of this corollary is additional.  The high-sector
proposition does not establish that premise, and the displayed time is not
a global Newman bound supplied by the proposition.

\begin{proposition}[Global patch relative to a published canopy]
Accept the proved canopy \(W(1/5)\le1/25\) in the proof of
\cite[Proposition~3.3 and Section~8]{polymath}, including the published
finite RH and barrier inputs used there.  Then
\[
 \Lambda\le\frac{11}{50}
      -\frac1{2500\cdot10^{32}+200}.
\]
\end{proposition}
\begin{proof}
The positive kernel gives \(H_t(iy)>0\), so a nonreal zero has nonzero
real part.  At a simple maximal zero \(z=x+iy\), with \(x,y>0\), its
distinct reflected pair \(-x\pm iy\) contributes exactly
\[
 E_t(z)\ge\frac1{2(x^2+w)},\qquad w=y^2.
\]
The same-height member contributes zero to \(E\); the other contributes
\(2/(4x^2+4w)\).  All other grouped contributions are nonnegative.
With \(X=10^{16}\), this supplies
\(E\ge1/[2(X^2+w)]\) on \(0<x\le X\), and the effective-probe
proposition supplies a stronger bound on \(x\ge X\).
Evenness covers negative real parts.  Combining by a spatial minimum
therefore gives the uniform floor
\[
 E_t(z)\ge\frac1{2(X^2+w)},\quad
 0<w\le1/25,\quad 1/5\le t\le3/10.
\]
We do not add the symmetry and probe estimates, since they can contain
the same roots.  Put \(a=X^2\) and
\[
 M_a(w)=\frac w4+\frac a8\log(1+2w/a),\qquad
 M_a'(w)=\frac{a+w}{2(a+2w)}.
\]
This derivative is the reciprocal speed corresponding to
\(w'\le-2-2w/(a+w)\).  The maximum-envelope comparison gives landing
by \(1/5+M_a(1/25)\), entirely within the time interval because
\(M_a(w)<w/2\).  The gain is conveniently bounded without subtraction
of nearly equal logarithms:
\[
 \frac w2-M_a(w)=\int_0^w\frac{u}{2(a+2u)}\,du
 \ge\frac{w^2}{4(a+2w)}.
\]
At \(w=1/25\), the last expression is \(1/(2500a+200)\), proving
the assertion.  Compactness and continuation use the positive-time
eventual-reality theorem and the maximum-envelope argument; no effective
value of the auxiliary compactness cutoff is needed.
\end{proof}

The gain in this historical comparison is about \(4\cdot10^{-36}\).
It is mathematically strict but quantitatively negligible.  We make no
record-bound or literature-novelty claim.  The actual earlier-time canopy
is the relevant published input: the bare statement \(\Lambda\le0.22\)
would not imply \(W(1/5)\le1/25\).  Increasing the high-sector field
above four alone does not improve the limiting low-sector floor.

\section{Two certificates on a compact positive-time interval}
\label{app:compact-certificates}

The following worked certificates illustrate the normalized value and
derivative alternatives.  They are local real-zero assertions.  They do
not cover all heights, exclude all complex zeros, or control times tending
to zero.  Both use \(1/5\le t\le3/10\), the symmetric normalization
\(Q_t=H_t/A_t\), and a fixed holomorphic approximant \(F_{t,N}\).

\begin{center}
\begin{tabular}{c|c|c|c}
Real interval & Fixed \(N\) & Natural cutoffs & Certified conclusion\\
\hline
\([451.9,452.5]\) & 6 & 5, 6 & \(Q_t>1/4\)\\
\([4005.2,4005.8]\) & 17 & 17 & \(Q'_t>2.3\)
\end{tabular}
\end{center}

For the first rectangle, outer disk center \(452.2\), radius \(R=.35\),
and inner radius \(a=.3\) give the full-disk error
\(\eta_6<2.248\).  The fixed approximant's value is bounded below by
\(2.507\) on 120 complete parameter cells (12 spatial by ten time cells).
Thus \(Q_t>2.507-2.248=.259>1/4\).  The natural cutoff changes across
\[
 x_6(t)=4\pi(36-t/16),
\]
whose entire range lies in
\([452.1537226679109918457,452.2322624842507366767]\).
The sixth summand and its reflected companion are paid in the full-disk
error wherever the natural cutoff is five.  No derivative of the moving
cutoff is taken.

For the second rectangle, outer center \(4005.5\), radius \(R=.7\),
and inner radius \(a=.3\) give \(\eta_{17}<.554\) and
\(|Q_t'-F_{t,17}'|<1.385\).  The entire outer rectangle has natural
cutoff 17.  Interval evaluation on 600 complete cells (30 spatial by
20 time cells) gives \(F_{t,17}'>3.70369\), hence \(Q_t'>2.3\).
The endpoints, enclosed over each of the 20 time cells, satisfy
\[
 Q_t(4005.2)<-.44,\qquad Q_t(4005.8)>.59.
\]
The intermediate value theorem and strict increase prove exactly one real
zero in \((4005.2,4005.8)\) for each time, and that zero is simple:
at a zero, \(H_t'=A_tQ_t'\ne0\).  Evenness gives both reflected
intervals.

Here are sufficient formulas to reproduce the finite evaluations.  For
real \(x>0\), put
\begin{align*}
 d&=1+x^2, & L&=\tfrac12\log d-\log(4\pi),
 & v&=\arctan(1/x),\\
 \alpha_r&=L/2-1/d,&
 \alpha_i&=-\pi/4+v/2+3x/d,\\
 \theta&=-7\pi/8-v/4+\tfrac x4(1-L)
                     +\tfrac t2\alpha_r\alpha_i,\\
 a_n&=\exp\{\tfrac t4\log^2n
                -(\tfrac12+\tfrac t2\alpha_r)\log n\},&
 \phi_n&=\theta+\tfrac12(x-t\alpha_i)\log n.
\end{align*}
Then \(F_{t,N}=2\sum_{n\le N}a_n\cos\phi_n\).
With
\[
 u=6(x^2-1)/d^2+1/d,\quad b=4x/d^2+x/d,\quad l_n=\log n,
\]
and \(L_n=(\alpha_r-l_n)(1+tu/2)-\alpha_i tb/2\), direct
differentiation gives the cancellation-aware expression
\[
 F_{t,N}'=\sum_{n\le N}a_n
        \left(L_n\sin\phi_n-\frac{tb l_n}{2}\cos\phi_n\right).
\]
These are exact consequences of the same logarithm branches used in the
analytic approximant; no finite-difference derivative is involved.

The certificates use 50-digit Decimal intervals with directed arithmetic
and widened correctly rounded elementary functions.  The software
guarantees for those elementary functions remain explicit computational
inputs.  A Machin identity supplies an enclosed value of \(\pi\), and
arctangent, sine, and cosine use Taylor series with explicit remainders.
The error bounds are established on the whole complex outer disk, including
the symmetric-normalization factor, reflection of the lower half-disk, and
every cutoff correction, before Cauchy's derivative bound is applied on
the inner interval.  Parameter-cell intervals cover the entire rectangles;
sampled values alone would not do so.

The regenerable certificates are
\href{../numerics/check_fixed_heat_checkpoint.py}{the cutoff-crossing checker}
and
\href{../numerics/fixed_heat_checkpoint_record_20261009.json}{its record},
and
\href{../numerics/check_zero_bearing_heat_cell.py}{the zero-bearing checker}
and
\href{../numerics/zero_bearing_heat_cell_record_20261009.json}{its record}.
The first script also contains unrelated fixed-scale bookkeeping; only
its heat certificate is used here.  Detailed parameters and outward
endpoints are retained in
\href{notes/4_CERTIFIED_COMPACT_CUTOFF_CROSSING_20261009.md}{Heat Note 4}
and
\href{notes/5_CERTIFIED_ZERO_BEARING_DERIVATIVE_CELL_20261009.md}{Heat Note 5}.
The imported approximation theorem and the analytic disk majorant are
mathematical inputs to these computations.  Internal replay and same-model
review do not constitute independent specialist validation of those inputs.


\section*{Acknowledgements and research provenance}
This manuscript was drafted for Edward Baker with substantial assistance
from GPT-6.1-sol (Codex), using reasoning effort ultra, for mathematical
derivation, source comparison, organization, and checking. Parallel
same-model reviews were used as internal checks and do not constitute
independent validation. The manuscript consolidates Heat Notes 1--12 and the dated research
notes in the \texttt{newman\_collisions} investigation and their scoped
reviews. Literature inputs are identified at the point of use; the
classical threshold theorem, the effective Polymath approximation,
the uniform count theorem, the local zero dynamics, and the cited
exponential-sum and Hilbert inequalities are not new
results of this manuscript. No claim of literature priority is made.
The source and small verification records are intended to remain
reviewable; a future independent mathematical review is still needed.

\begin{thebibliography}{99}
\bibitem{polymath}
D. H. J. Polymath,
\emph{Effective approximation of heat flow evolution of the Riemann
$\xi$ function, and a new upper bound for the de Bruijn--Newman constant},
2019, arXiv:1904.12438v2.
\url{https://arxiv.org/abs/1904.12438v2}.

\bibitem{rodgerstao}
B. Rodgers and T. Tao,
\emph{The de Bruijn--Newman constant is non-negative},
2018, arXiv:1801.05914.
\url{https://arxiv.org/abs/1801.05914}.

\bibitem{newmanwu}
C. M. Newman and W. Wu,
\emph{Constants of de Bruijn--Newman type in analytic number theory
and statistical physics}, 2019, arXiv:1901.06596.
\url{https://arxiv.org/abs/1901.06596}.

\bibitem{kikimlee}
H. Ki, Y.-O. Kim, and J. Lee,
\emph{On the de Bruijn--Newman constant},
Advances in Mathematics \textbf{222} (2009), 281--306.
\url{https://doi.org/10.1016/j.aim.2009.04.003}.

\bibitem{gasper}
G. Gasper,
\emph{Using integrals of squares of certain real-valued special functions
to prove that the P\'olya $\Xi^*(z)$ function, the functions $K_{iz}(a)$,
$a>0$, and some other entire functions have only real zeros},
2008, arXiv:0801.2996.
\url{https://arxiv.org/abs/0801.2996}.

\bibitem{planat}
M. Planat,
\emph{Collective Contraction in the de Bruijn--Newman Heat Flow:
Off-Zero Logarithmic Derivatives and Certified Barrier Refinements},
2026, arXiv:2609.37164v2 (preprint).
\url{https://arxiv.org/abs/2609.37164v2}.
\bibitem{arias}
J. Arias de Reyna,
\emph{Explicit van der Corput's $d$-th derivative estimate},
2024, arXiv:2407.02094v1.
\url{https://arxiv.org/html/2407.02094v1}.

\bibitem{trudgianyang}
T. S. Trudgian and A. Yang,
\emph{Toward optimal exponent pairs}, arXiv:2306.05599v3,
Sections 1.1, 1.2, and 1.4.
\url{https://arxiv.org/html/2306.05599v3}.

\bibitem{montgomeryvaughan}
H. L. Montgomery and R. C. Vaughan,
\emph{Hilbert's Inequality},
Journal of the London Mathematical Society (2) \textbf{8} (1974), 73--82.
\url{https://personal.science.psu.edu/rcv4/personal/Publications/s2-8-1-73.pdf}.

\bibitem{dlmf}
NIST Digital Library of Mathematical Functions,
\emph{Asymptotic Formulas: Primes}, Section 27.12, equation 27.12.4
(prime number theorem, leading term).
\url{https://dlmf.nist.gov/27.12\#E4}.

\end{thebibliography}
\end{document}
